\documentclass[11pt]{amsart}

\usepackage[T1]{fontenc}
\usepackage[utf8]{inputenc}
\usepackage{lmodern}
\usepackage[a4paper,margin=1.05in]{geometry}
\usepackage{microtype}
\usepackage{amsmath,amssymb,amsthm,mathtools,mathrsfs}
\usepackage{enumitem}
\usepackage{booktabs}
\usepackage[dvipsnames]{xcolor}
\usepackage[colorlinks=true,linkcolor=MidnightBlue,citecolor=ForestGreen,urlcolor=BrickRed]{hyperref}
\usepackage{array}
\usepackage{longtable}
\usepackage{tikz-cd}

\numberwithin{equation}{section}

\newtheorem{theorem}{Theorem}[section]
\newtheorem{proposition}[theorem]{Proposition}
\newtheorem{corollary}[theorem]{Corollary}
\newtheorem{lemma}[theorem]{Lemma}
\newtheorem{conjecture}[theorem]{Conjecture}
\theoremstyle{definition}
\newtheorem{definition}[theorem]{Definition}
\newtheorem{example}[theorem]{Example}
\newtheorem{remark}[theorem]{Remark}
\newtheorem{warning}[theorem]{Warning}

\newcommand{\B}{\mathbf B}
\newcommand{\Z}{\mathbf Z}
\newcommand{\Q}{\mathbf Q}
\newcommand{\R}{\mathbf R}
\newcommand{\C}{\mathbf C}
\newcommand{\N}{\mathbf N}
\newcommand{\F}{\mathbf F}
\newcommand{\No}{\mathbf{No}}
\newcommand{\Spec}{\operatorname{Spec}}
\newcommand{\Proj}{\operatorname{Proj}}
\newcommand{\supp}{\operatorname{supp}}
\newcommand{\Supp}{\operatorname{Supp}}
\newcommand{\im}{\operatorname{Im}}
\newcommand{\kerp}{\operatorname{Ker}}
\newcommand{\Hom}{\operatorname{Hom}}
\newcommand{\End}{\operatorname{End}}
\newcommand{\Aut}{\operatorname{Aut}}
\newcommand{\GL}{\operatorname{GL}}
\newcommand{\PGL}{\operatorname{PGL}}
\newcommand{\Frac}{\operatorname{Frac}}
\newcommand{\ord}{\operatorname{ord}}
\newcommand{\st}{\operatorname{st}}
\newcommand{\gr}{\operatorname{gr}}
\newcommand{\cO}{\mathcal O}
\newcommand{\cM}{\mathcal M}
\newcommand{\cU}{\mathcal U}
\newcommand{\zero}{\mathbf 0}
\newcommand{\Pum}{\mathbf P_{\mathrm{um}}}
\newcommand{\Sp}{\operatorname{Sp}}
\newcommand{\id}{\operatorname{id}}
\newcommand{\colim}{\operatorname*{colim}}
\newcommand{\liminv}{\operatorname*{lim}}
\newcommand{\eps}{\varepsilon}
\newcommand{\taup}{\tau}
\newcommand{\eR}{e}
\newcommand{\ideal}[1]{\mathfrak #1}

\DeclareMathOperator{\Bool}{Bool}
\DeclareMathOperator{\Mon}{Mon}
\DeclareMathOperator{\Ring}{Ring}
\DeclareMathOperator{\SemiRing}{SemiRing}

\title[Split-Zero Projective Geometry and Leibniz Monads]
{Split-Zero Projective Geometry and Leibniz Monads\\[4pt]
\large Boolean-Cubical Localization, \texorpdfstring{$\mathbb F_1$}{F1}-Skeletons, and Surcomplex Infinitesimal Fibres}
\author{The Clankers}
\date{June 2026\\Pre-publication draft}

\begin{document}

\begin{abstract}
Let $R$ be a nonzero commutative ring and let
\[
G(R)=R\sqcup\{\tau\}
\]
be its split-zero globalization, in which $\tau$ is the semiring zero and the ring zero $e=0_R$ remains as a distinct supported element.  Earlier work established the ideal theory, prime spectrum, ring reflection, Boolean support character, and the classification of $G(R)$-semimodules by join-semilattice-indexed diagrams of $R$-modules.  The present article develops the geometric and infinitesimal consequences that are not contained in those notes.

The first main result is the exact division trichotomy.  Inverting a classically nonzero scalar gives ordinary ring localization, inverting the supported zero gives the Boolean semiring, and inverting the actual semiring zero gives the one-element semiring.  In finite products this becomes a coordinatewise mixed-support localization theorem, and adjoining an identity between a supported zero mask and a nonzero amplitude Booleanizes precisely the selected face of the Boolean cube.

We then construct a support-stratified projective geometry.  For a field $K$ the unimodular projective point set decomposes canonically as
\[
\Pum^n(G(K))\cong
\bigsqcup_{\varnothing\neq A\subseteq\{0,\dots,n\}}
\mathbf P^{|A|-1}(K).
\]
Over $\F_q$ this yields an explicit point-counting polynomial and a support-projective zeta function.  The linear automorphism group is exactly the monomial group
\[
\GL_{n+1}(G(K))\cong (K^\times)^{n+1}\rtimes S_{n+1},
\]
so the global split projective symmetry is the toric or $\mathbb F_1$ symmetry, rather than the full classical projective linear group.

Finally, for a residue map $A\twoheadrightarrow k$ with infinitesimal ideal $\mathfrak m$, we prove that the Leibniz monad is the supported fibre $a+\mathfrak m$, while $\tau$ records absence and the Boolean cube records possible activation of absent coordinates.  Formal completion commutes with split-zero globalization.  Applied to finite hyperreals and to the finite surcomplex valuation ring $\cO_{\No[i]}$, this produces Boolean cubes of ordinary infinitesimal monads.  In the surcomplex case each nonzero infinitesimal has an order and complex phase, whereas the supported zero is the valuation-infinite centre of the monad.  Projective inversion identifies the pointed zero and infinity monads, and under univalence this equivalence becomes an identity of structured objects without identifying their arithmetic centres.
\end{abstract}

\maketitle
\tableofcontents

\section{Introduction and statement of the results}

The split-zero semiring
\[
G(R)=R\sqcup\{\tau\}
\]
was introduced and developed in the globalization notes \cite{Globalization,Secondary}.  The defining feature is the strict distinction
\[
\tau=0_{G(R)}\neq e:=0_R\in R\subset G(R).
\]
The element $\tau$ denotes absence of support, while $e$ denotes an arithmetically zero value carried by nonempty support.  The baseline note \cite{Globalization} classifies ideals, prime ideals, congruences, and localizations for $G(\Z)$ and its number-field variants.  The secondary note \cite{Secondary} proves, among other results, that the ring reflection
\[
p_R:G(R)\longrightarrow R
\]
is universal among maps to rings, that the support character
\[
\chi_R:G(R)\longrightarrow\B
\]
is the unique Boolean character, and that $G(R)$-semimodules are exactly diagrams of $R$-modules over join-semilattices.  In particular, the free semimodule $G(R)^n$ has support poset the Boolean cube $\B^n$.

The purpose of the present paper is to develop the geometry forced by those algebraic facts.  The results are not a reformulation of ordinary ring geometry.  The two zero states produce two different principal opens, two different projective poles, two different notions of vanishing, and a cubical infinitesimal neighbourhood theory.

The results proved below may be summarized as follows.

\begin{enumerate}[label=\textup{(\arabic*)},leftmargin=3.2em]
\item \textbf{Division trichotomy.}
For every nonzero ring $R$,
\[
G(R)[e^{-1}]\cong\B,
\qquad
G(R)[\tau^{-1}]\cong\zero.
\]
If $M\subset R$ is multiplicative and $0\notin M$, then
\[
G(R)[M^{-1}]\cong G(M^{-1}R).
\]
Thus ordinary nonzero localization preserves amplitudes, division by $e$ retains only support, and division by $\tau$ destroys every distinction.

\item \textbf{Absolute infinitesimals.}
The kernel of the universal ring reflection is
\[
\{\tau,e\}.
\]
In $G(R)^I$ it is the Boolean cube
\[
\{\eps_A:A\subseteq I\}\cong\B^I.
\]
This identifies $e$ as the unique scalar element that is nonzero in the semiring but zero in every ring-valued realization.  It is an idempotent support residue, not a tangent displacement.

\item \textbf{Mixed-support localization and scoped identity.}
For a field $K$, $S=G(K)$, a finite set $I$, and $x\in S^I$, let $U(x)$ be the coordinates at which $x$ has nonzero amplitude, $Z(x)$ the coordinates equal to $e$, and $N(x)$ the coordinates equal to $\tau$.  Then
\[
S^I[x^{-1}]
\cong
S^{U(x)}\times\B^{Z(x)}.
\]
Coordinates in $N(x)$ become terminal and disappear.  If an equality $e=h_i$ with $h_i\in K^\times$ is imposed only on a subset $A\subseteq I$, the quotient is
\[
S^I/(\eps_A=h_A)
\cong
\B^A\times S^{I\setminus A}.
\]
Thus an identity may be adjoined locally on a face without collapsing the complementary geometry.

\item \textbf{Generic-point suspension.}
For every nonzero ring $R$,
\[
\Spec G(R)=\{P_\tau\}\sqcup\{P_{\mathfrak p}:\mathfrak p\in\Spec R\},
\]
where $P_\tau=\{\tau\}$ is an open dense generic point and $P_{\mathfrak p}=\mathfrak p\cup\{\tau\}$.  The closed subspace $V(e)$ is canonically $\Spec R$, while
\[
D(e)=\{P_\tau\}.
\]
For $R=\Z$, the classical generic point $(0)$ becomes the nonzero prime $\{\tau,0\}$, and $\{\tau\}$ becomes the new zero ideal.

\item \textbf{Projective Boolean-face decomposition.}
For every field $K$,
\[
\Pum^n(G(K))
\cong
\bigsqcup_{\varnothing\neq A\subseteq[n]_0}
\mathbf P^{|A|-1}(K),
\qquad [n]_0:=\{0,\dots,n\}.
\]
The support map to $\Pum^n(\B)$ has these projective spaces as its fibres.

\item \textbf{Point counts and zeta.}
For $K=\F_q$,
\[
\#\Pum^n(G(\F_q))
=
\frac{(q+1)^{n+1}-2^{n+1}}{q-1}.
\]
The associated support-projective zeta function is
\[
Z_{\mathrm{sp}}(\Pum^n/\F_q,T)
=
\prod_{j=0}^{n}(1-q^jT)^{-c_{n,j}},
\qquad
c_{n,j}=\sum_{r=j+1}^{n+1}\binom{n+1}{r}.
\]

\item \textbf{Toric linear symmetry.}
Every invertible matrix over $G(K)$ is monomial:
\[
\GL_m(G(K))\cong(K^\times)^m\rtimes S_m.
\]
Consequently the global split-projective linear group is the normalizer of the diagonal torus.  General translations, general M\"obius maps, and the Cayley transform act on the full-support classical stratum but do not extend to global linear automorphisms of the mixed-support completion.

\item \textbf{Leibniz monads.}
For a surjection $q:A\twoheadrightarrow k$ with kernel $\mathfrak m$, the supported fibre over $a\in k$ is
\[
\mu(a)=\widetilde a+\mathfrak m.
\]
For a point with support $J\subseteq I$, the upward mixed-support monad is a Boolean cube of such fibres:
\[
\mu^\square(a)
\cong
\bigsqcup_{L\subseteq I\setminus J}
\left(\prod_{i\in J}(\widetilde a_i+\mathfrak m)\right)
\times \mathfrak m^L\times\{\tau\}^{I\setminus(J\cup L)}.
\]
Thus $e$ is the centre of a monad, $\mathfrak m$ is its infinitesimal halo, and $\tau$ is the absent face.

\item \textbf{Surcomplex order and phase.}
For $K=\No[i]$, every nonzero element has a Conway--Hahn expansion
\[
z=\sum_{\rho<\lambda}c_\rho\omega^{\gamma_\rho},
\qquad c_\rho\in\C^\times,
\qquad \gamma_0>\gamma_1>\cdots.
\]
The natural valuation gives a residue map
\[
\st:\cO_K\twoheadrightarrow\C
\]
whose kernel is the ideal of surcomplex infinitesimals.  Each valuation layer is one-dimensional over $\C$, and reciprocal ordinals are coinitial among positive surreal infinitesimals.  Split-zero globalization adds an independent Boolean support coordinate to this order-and-phase arithmetic.

\item \textbf{Zero, infinity, and univalence.}
Projective inversion exchanges the supported and unsupported zero points with corresponding supported and unsupported infinity points.  It yields an isomorphism of pointed zero and infinity monads.  In a univalent universe, that equivalence gives an identity of the structured monads, but it does not identify the arithmetic elements $e$ and $\omega^{-1}$; imposing that arithmetic equality produces the Boolean quotient.
\end{enumerate}

All new claims are proved below.  Standard facts about surreal normal forms and the Berarducci--Mantova derivation are imported from \cite{Conway,BerarducciMantova}; standard facts about univalence are imported from \cite{HoTT}.  No result concerning surreal numbers is used in the proofs of the split-zero algebraic or projective theorems.

\section{Split-zero algebra and universal maps}

Throughout, a semiring is commutative, unital, and has a specified zero; homomorphisms preserve both $0$ and $1$.  The one-element semiring, in which $0=1$, is denoted by $\zero$ and is allowed.

\begin{definition}[Split-zero globalization]
Let $R$ be a nonzero commutative ring.  Define
\[
S:=G(R):=R\sqcup\{\tau\},
\qquad \tau\notin R,
\]
with operations
\[
r\oplus s=r+s,
\qquad
r\oplus\tau=\tau\oplus r=r,
\qquad
\tau\oplus\tau=\tau,
\]
and
\[
rs=rs,
\qquad
r\tau=\tau r=\tau,
\qquad
\tau^2=\tau.
\]
We write
\[
e:=0_R\in R\subset S.
\]
Thus $0_S=\tau$ and $e\neq\tau$.
\end{definition}

\begin{proposition}[Pair normal form]\label{prop:pair-normal-form}
There is a canonical semiring isomorphism
\[
G(R)\xrightarrow{\sim}
D_R:=\{(0,0)\}\cup(R\times\{1\})
\subseteq R\times\B
\]
given by
\[
\tau\longmapsto(0,0),
\qquad
r\longmapsto(r,1).
\]
Under this isomorphism,
\[
e\longmapsto(0,1).
\]
\end{proposition}

\begin{proof}
The set $D_R$ is closed under the componentwise operations of $R\times\B$.  The displayed map is bijective, and the transported operations are exactly the defining operations of $G(R)$.
\end{proof}

\begin{definition}[Reflection and support]
The ring reflection and support character are
\[
p_R:G(R)\longrightarrow R,
\qquad
p_R(\tau)=0,
\quad
p_R(r)=r,
\]
and
\[
\chi_R:G(R)\longrightarrow\B,
\qquad
\chi_R(\tau)=0,
\quad
\chi_R(r)=1.
\]
\end{definition}

The following two universal properties were proved in \cite{Secondary}; they are recorded because they control every later construction.

\begin{theorem}[Universal ring reflection and Boolean character]\label{thm:universal-maps}
Let $R$ be a nonzero commutative ring.
\begin{enumerate}[label=\textup{(\alph*)}]
\item For every commutative ring $A$, composition with $p_R$ induces a bijection
\[
\Hom_{\Ring}(R,A)
\xrightarrow{\sim}
\Hom_{\SemiRing}(G(R),A).
\]
\item The support map $\chi_R$ is the unique semiring homomorphism $G(R)\to\B$.
\end{enumerate}
\end{theorem}

\begin{proof}
Part (a) is Theorem 2.2 of \cite{Secondary}.  For completeness, if $f:G(R)\to A$ has ring target, then
\[
f(e)=f(e\oplus e)=f(e)+f(e),
\]
so subtraction gives $f(e)=0_A$.  Hence $f|_R$ is a ring homomorphism and $f=(f|_R)\circ p_R$.

For part (b), let $u:G(R)\to\B$.  Since $1\oplus(-1)=e$ in $G(R)$,
\[
u(e)=u(1)+u(-1)=1+u(-1)=1.
\]
If $u(r)=0$ for some supported $r\in R$, then
\[
1=u(e)=u(re)=u(r)u(e)=0,
\]
a contradiction.  Hence every supported element maps to $1$, while $\tau$ maps to $0$.
\end{proof}

\section{Idempotent localization and the division trichotomy}

We use localization only through its universal property.  If $A$ is a semiring and $M\subseteq A$ is multiplicatively closed, a localization of $A$ at $M$ is a semiring $M^{-1}A$ with a map $\lambda:A\to M^{-1}A$ such that every $\lambda(m)$ is invertible and every map from $A$ that inverts $M$ factors uniquely through $\lambda$.

\begin{theorem}[Localization at an idempotent]\label{thm:idempotent-localization}
Let $A$ be a commutative semiring and let $a\in A$ satisfy $a^2=a$.  Then
\[
A[a^{-1}]\cong aA,
\]
where $aA=\{ax:x\in A\}$ is regarded as a semiring with zero $a0_A$ and multiplicative identity $a$.
\end{theorem}

\begin{proof}
The subset $aA$ is closed under addition and multiplication because
\[
ax+ay=a(x+y),
\qquad
(ax)(ay)=a^2xy=axy.
\]
The element $a$ is its multiplicative identity.  Define
\[
\pi_a:A\longrightarrow aA,
\qquad
x\longmapsto ax.
\]
Then
\[
\pi_a(a)=a^2=a=1_{aA},
\]
so $a$ becomes invertible.

Let $f:A\to T$ be a semiring homomorphism such that $f(a)$ is invertible.  Since $f(a)$ is idempotent,
\[
f(a)^2=f(a).
\]
Multiplication by $f(a)^{-1}$ gives $f(a)=1_T$.  Define
\[
\overline f:aA\longrightarrow T,
\qquad
\overline f(ax)=f(x).
\]
If $ax=ay$, then
\[
f(a)f(x)=f(a)f(y),
\]
and $f(a)=1_T$, so $f(x)=f(y)$.  Thus $\overline f$ is well defined.  It is a semiring homomorphism and satisfies $f=\overline f\circ\pi_a$.  Uniqueness follows because every element of $aA$ has the form $\pi_a(x)$.
\end{proof}

\begin{corollary}[The two zero localizations]\label{cor:two-zero-localizations}
For every nonzero commutative ring $R$,
\[
G(R)[e^{-1}]\cong\{\tau,e\}\cong\B,
\]
and
\[
G(R)[\tau^{-1}]\cong\{\tau\}\cong\zero.
\]
Under the first isomorphism, the localization map is the support character:
\[
\frac{x}{e}=\chi_R(x).
\]
\end{corollary}

\begin{proof}
Both $e$ and $\tau$ are idempotent.  Moreover,
\[
eG(R)=\{\tau,e\},
\qquad
\tau G(R)=\{\tau\}.
\]
In $eG(R)$ the element $e$ is the identity and $\tau$ is the zero, so $eG(R)\cong\B$.  The localization map is $x\mapsto ex$, which sends $\tau$ to $\tau$ and every supported element to $e$; this is exactly $\chi_R$ after the identification $\tau\leftrightarrow0$, $e\leftrightarrow1$.
\end{proof}

\begin{theorem}[Compatibility with ordinary ring localization]\label{thm:ring-localization}
Let $M\subseteq R$ be a multiplicatively closed set with $0\notin M$.  Then the canonical map $R\to M^{-1}R$ induces a canonical semiring isomorphism
\[
M^{-1}G(R)\xrightarrow{\sim}G(M^{-1}R).
\]
\end{theorem}

\begin{proof}
Let
\[
\lambda:G(R)\longrightarrow G(M^{-1}R)
\]
be defined by $\lambda(\tau)=\tau$ and $\lambda(r)=r/1$.  Every $m\in M$ maps to the supported unit $m/1$.

Let $f:G(R)\to T$ be a semiring homomorphism such that every $f(m)$ is invertible.  Define
\[
\overline f(\tau):=0_T,
\qquad
\overline f(r/m):=f(r)f(m)^{-1}.
\]
If $r/m=s/n$ in $M^{-1}R$, then there exists $u\in M$ with
\[
u(nr-ms)=0
\]
in $R$, equivalently $unr=ums$.  Applying $f$ gives
\[
f(u)f(n)f(r)=f(u)f(m)f(s).
\]
Since $f(u)$, $f(m)$, and $f(n)$ are invertible, one obtains
\[
f(r)f(m)^{-1}=f(s)f(n)^{-1}.
\]
Thus $\overline f$ is well defined.  The usual fraction calculations show that it is a semiring homomorphism.  It is the unique map satisfying $\overline f\circ\lambda=f$, because the supported fractions and $\tau$ exhaust $G(M^{-1}R)$.
\end{proof}

\begin{corollary}[Classical fraction field versus Boolean fraction residue]\label{cor:fraction-comparison}
Let $R$ be an integral domain with fraction field $K$.
\begin{enumerate}[label=\textup{(\alph*)}]
\item Localizing at the classically nonzero elements gives
\[
G(R)[(R\setminus\{0\})^{-1}]
\cong
G(K).
\]
\item Localizing at all semiring-nonzero elements, namely at $R=G(R)\setminus\{\tau\}$, gives
\[
G(R)[R^{-1}]
\cong
\B.
\]
\end{enumerate}
\end{corollary}

\begin{proof}
Part (a) is Theorem \ref{thm:ring-localization}.  For part (b), the denominator set contains $e$, so the localization factors through $G(R)[e^{-1}]\cong\B$.  In $\B$, every supported element of $G(R)$ already maps to $1$, hence is invertible.  Therefore $\B$ has the universal property of the localization at all supported elements.
\end{proof}

\begin{remark}[Element localization versus localization at a prime]\label{rem:prime-v-element}
For $R=\Z$, let
\[
P_0=\{\tau,0\}.
\]
The stalk localization at $P_0$ inverts the complement $\Z\setminus\{0\}$ and is therefore $G(\Q)$.  Principal localization at the element $e=0$ instead gives $\B$.  Thus ``localizing at the generic zero prime'' and ``placing the zero element in the denominator set'' are distinct operations.
\end{remark}

\begin{theorem}[Strict identification with a nonzero amplitude]\label{thm:strict-identification-collapse}
Let $K$ be a field, let $S=G(K)$, and let $h\in K^\times$.  Let $\sim$ be the smallest semiring congruence on $S$ satisfying
\[
e\sim h.
\]
Then
\[
S/{\sim}\cong\B.
\]
\end{theorem}

\begin{proof}
In the quotient, the class of $h$ is invertible, so the class of $e$ is invertible.  Since $e^2=e$, its class is an invertible idempotent and therefore equals $1$.  For every supported $x\in K$,
\[
ex=e
\]
in $S$, hence in the quotient
\[
x=1x=ex=e=1.
\]
Thus all supported elements lie in one congruence class.  The support character $\chi_K:S\to\B$ satisfies $\chi_K(e)=\chi_K(h)=1$, so it coequalizes the imposed relation and distinguishes $\tau$ from the supported class.  Hence the quotient has exactly two elements and is Boolean.
\end{proof}

\begin{theorem}[Division trichotomy]\label{thm:division-trichotomy}
Let $R$ be an integral domain with fraction field $K$.  The three universal division operations are
\[
\boxed{
\begin{aligned}
\text{invert }R\setminus\{0\}
&\quad\Longrightarrow\quad G(K),\\
\text{invert }e=0_R
&\quad\Longrightarrow\quad \B,\\
\text{invert }\tau=0_{G(R)}
&\quad\Longrightarrow\quad \zero.
\end{aligned}}
\]
The first preserves amplitude and support, the second preserves only support, and the third preserves no distinction.
\end{theorem}

\begin{proof}
This is the conjunction of Corollaries \ref{cor:two-zero-localizations} and \ref{cor:fraction-comparison}.
\end{proof}

\section{Ring-null infinitesimals and Boolean cubes}

The preceding localization theorem suggests an intrinsic infinitesimal notion which is independent of order, topology, or valuation.

\begin{definition}[Absolute ring-null halo]\label{def:absolute-halo}
Let $A$ be a commutative semiring admitting a universal ring reflection
\[
\rho_A:A\longrightarrow A^{\mathrm{ring}}.
\]
Define
\[
\mathfrak I_{\mathrm{abs}}(A):=\rho_A^{-1}(0).
\]
An element of $\mathfrak I_{\mathrm{abs}}(A)\setminus\{0_A\}$ is called an \emph{absolute ring-null infinitesimal}.
\end{definition}

\begin{proposition}[Intrinsic characterization]\label{prop:absolute-characterization}
For $x\in A$,
\[
x\in\mathfrak I_{\mathrm{abs}}(A)
\]
if and only if every semiring homomorphism from $A$ to a commutative ring sends $x$ to zero.
\end{proposition}

\begin{proof}
Every map from $A$ to a ring factors uniquely through $\rho_A$.  Hence $\rho_A(x)=0$ implies that every ring-valued map kills $x$.  Conversely, apply the assumed condition to the ring-valued map $\rho_A$ itself.
\end{proof}

\begin{corollary}[The scalar split-zero infinitesimal]\label{cor:scalar-absolute-halo}
For $S=G(R)$,
\[
\mathfrak I_{\mathrm{abs}}(S)=\{\tau,e\}.
\]
Thus $e$ is the unique nonzero absolute ring-null infinitesimal of $S$.
\end{corollary}

\begin{proof}
The universal ring reflection is $p_R$ by Theorem \ref{thm:universal-maps}, and
\[
p_R^{-1}(0_R)=\{\tau,e\}.
\]
The semiring zero is $\tau$, so $e$ is the unique nonzero member of this fibre.
\end{proof}

Let $I$ be a finite set and write
\[
S^I=\prod_{i\in I}S.
\]
For $A\subseteq I$, define the supported-zero mask
\[
(\eps_A)_i=
\begin{cases}
e,&i\in A,\\
\tau,&i\notin A.
\end{cases}
\]

\begin{theorem}[Finite-product ring reflection]\label{thm:product-ring-reflection}
The coordinatewise map
\[
p_I:=p_R^I:S^I\longrightarrow R^I
\]
is the universal ring reflection of $S^I$.  Its kernel is
\[
p_I^{-1}(0)=\{\eps_A:A\subseteq I\}.
\]
Moreover,
\[
\eps_A+\eps_B=\eps_{A\cup B},
\qquad
\eps_A\eps_B=\eps_{A\cap B}.
\]
Consequently,
\[
\mathfrak I_{\mathrm{abs}}(S^I)
\cong\B^I
\]
as semirings.
\end{theorem}

\begin{proof}
Let $f:S^I\to T$ be a semiring homomorphism with ring target.  Each $\eps_A$ is additively idempotent:
\[
\eps_A+\eps_A=\eps_A.
\]
Therefore
\[
f(\eps_A)+f(\eps_A)=f(\eps_A),
\]
and subtraction in $T$ gives $f(\eps_A)=0$.

For $r=(r_i)_{i\in I}\in R^I$, let $\widetilde r\in S^I$ be its fully supported lift, so $(\widetilde r)_i=r_i\in R\subset S$.  Define
\[
g:R^I\longrightarrow T,
\qquad
g(r):=f(\widetilde r).
\]
Because sums and products of fully supported lifts remain fully supported,
\[
\widetilde{r+s}=\widetilde r+\widetilde s,
\qquad
\widetilde{rs}=\widetilde r\,\widetilde s,
\]
where the operations on $R^I$ are coordinatewise.  Hence $g$ is a ring homomorphism.

If $x\in S^I$ and $r=p_I(x)$, then $x$ and $\widetilde r$ differ only in coordinates where $r_i=0$: in such a coordinate $x_i$ is either $\tau$ or $e$, while $(\widetilde r)_i=e$.  Therefore there is a mask $\eps_A$ such that
\[
\widetilde r=x+\eps_A.
\]
Applying $f$ gives
\[
f(\widetilde r)=f(x)+f(\eps_A)=f(x).
\]
Thus $f=g\circ p_I$.  Uniqueness follows from surjectivity of $p_I$.

The kernel description is immediate coordinatewise.  The union and intersection formulas follow from
\[
e+e=e,
\quad e+\tau=e,
\quad \tau+\tau=\tau,
\]
and
\[
e^2=e,
\quad e\tau=\tau,
\quad \tau^2=\tau.
\]
These are exactly the join and meet operations on $\B^I$.
\end{proof}

\begin{example}[The mixed zero $(\tau,e)$]
In $S^2$ the four absolute infinitesimals are
\[
(\tau,\tau),
\quad
(e,\tau),
\quad
(\tau,e),
\quad
(e,e).
\]
The element $(\tau,e)$ means that the first coordinate is absent while the second coordinate is supported but arithmetically zero.  For $a\in R$,
\[
(a,\tau)+(-a,\tau)=(e,\tau),
\]
so cancellation produces supported zero rather than absence.
\end{example}

\subsection{Derivations and the failure of tangent displacement}

\begin{definition}[Semiring derivation]
Let $A$ be a semiring and let $M$ be an $A$-semimodule.  A derivation is a map $d:A\to M$ satisfying
\[
d(0)=0,
\qquad
d(1)=0,
\qquad
d(x+y)=d(x)+d(y),
\]
and
\[
d(xy)=x\,d(y)+y\,d(x).
\]
\end{definition}

\begin{proposition}[The supported zero is derivationally constant]\label{prop:derivation-kills-e}
Every derivation $d:G(R)\to M$ satisfies
\[
d(e)=0.
\]
Every derivation $\B\to M$ is zero.
\end{proposition}

\begin{proof}
Since $1+e=1$ in $G(R)$,
\[
0=d(1)=d(1+e)=d(1)+d(e)=d(e).
\]
For $\B$, the only elements are $0$ and $1$, both of which are killed by the defining axioms of a derivation.
\end{proof}

\begin{remark}
The identity $x+e=x$ for every supported $x\in R$ shows that $e$ does not move a supported point.  It is therefore an infinitesimal only in the ring-null support sense of Definition \ref{def:absolute-halo}; it is not a first-order tangent parameter.
\end{remark}

The closest ring-valued multiplicative model of the relation $e^2=e$ makes this distinction exact.

\begin{theorem}[Idempotent probes are finite-difference probes]\label{thm:idempotent-probe}
Let $A$ be a commutative ring and set
\[
C:=A[q]/(q^2-q).
\]
Then evaluation at $q=0$ and $q=1$ induces a canonical ring isomorphism
\[
C\xrightarrow{\sim}A\times A.
\]
For every polynomial $f\in A[T]$ and every $x,h\in A$,
\[
f(x+qh)
=
f(x)+q\bigl(f(x+h)-f(x)\bigr)
\]
in $C$.  Moreover,
\[
\Omega^1_{C/A}=0.
\]
\end{theorem}

\begin{proof}
The ideals $(q)$ and $(q-1)$ of $A[q]$ are comaximal because
\[
q-(q-1)=1.
\]
Their intersection equals their product $(q(q-1))=(q^2-q)$.  The Chinese remainder theorem gives
\[
A[q]/(q^2-q)
\cong
A[q]/(q)\times A[q]/(q-1)
\cong
A\times A.
\]
Under this isomorphism, $q\mapsto(0,1)$ and $x+qh\mapsto(x,x+h)$.  Both sides of the polynomial identity map to
\[
\bigl(f(x),f(x+h)\bigr),
\]
so they are equal.

For relative K\"ahler differentials, the relation $q^2-q=0$ gives
\[
(2q-1)dq=0.
\]
But
\[
(2q-1)^2=4q^2-4q+1=1
\]
in $C$, so $2q-1$ is invertible and $dq=0$.  Since $C$ is generated over $A$ by $q$, this implies $\Omega^1_{C/A}=0$.
\end{proof}

\begin{corollary}[Idempotence versus nilpotence]
An idempotent parameter records two endpoint values and exact finite differences.  A nilpotent parameter $\epsilon^2=0$ records first-order tangent data.  These structures cannot be identified in a nonzero ring.
\end{corollary}

\begin{proof}
The finite-difference statement is Theorem \ref{thm:idempotent-probe}.  If a ring element satisfies both $q^2=q$ and $q^2=0$, then $q=0$.
\end{proof}

\section{Mixed-support localization and scoped identities}

\begin{lemma}[Localization of a finite product]\label{lem:product-localization}
Let $A_i$ be commutative semirings indexed by a finite set $I$, let
\[
A:=\prod_{i\in I}A_i,
\qquad
x=(x_i)_{i\in I}\in A.
\]
Then
\[
A[x^{-1}]
\cong
\prod_{i\in I}A_i[x_i^{-1}].
\]
\end{lemma}

\begin{proof}
Use the standard fraction construction for localization.  Every element of $A[x^{-1}]$ is represented by a fraction
\[
\frac{y}{x^n},
\qquad y\in A,
\quad n\ge0.
\]
Define
\[
\Phi\left(\frac{y}{x^n}\right)
:=
\left(\frac{y_i}{x_i^n}\right)_{i\in I}.
\]
If two fractions are equal in $A[x^{-1}]$, their coordinate fractions are equal, so $\Phi$ is well defined.

For surjectivity, choose fractions $y_i/x_i^{n_i}$ in the factors and put $N=\max_i n_i$.  Replacing $y_i$ by $y_i x_i^{N-n_i}$ gives a common exponent $N$, hence a preimage in $A[x^{-1}]$.

For injectivity, suppose the coordinate fractions of $y/x^n$ and $z/x^m$ are equal.  For each $i$ there exists $k_i\ge0$ such that
\[
x_i^{k_i+m}y_i=x_i^{k_i+n}z_i.
\]
Let $k=\max_i k_i$.  Multiplying each equality by a further power of $x_i$ yields
\[
x_i^{k+m}y_i=x_i^{k+n}z_i
\]
for every $i$.  Therefore
\[
x^{k+m}y=x^{k+n}z
\]
in $A$, which is precisely the equality of the original fractions in $A[x^{-1}]$.
\end{proof}

Let $K$ be a field, $S=G(K)$, and $I$ finite.  For $x\in S^I$, define
\[
U(x):=\{i\in I:x_i\in K^\times\},
\]
\[
Z(x):=\{i\in I:x_i=e\},
\]
and
\[
N(x):=\{i\in I:x_i=\tau\}.
\]
These form a partition of $I$.

\begin{theorem}[Mixed-support localization normal form]\label{thm:mixed-localization}
For every $x\in S^I$,
\[
S^I[x^{-1}]
\cong
S^{U(x)}\times\B^{Z(x)}.
\]
The coordinates in $N(x)$ contribute one-element factors and hence disappear.  Under the displayed isomorphism,
\[
\frac{y}{x}
\longmapsto
\left(
\left(\frac{y_i}{x_i}\right)_{i\in U(x)},
\left(\chi_K(y_i)\right)_{i\in Z(x)}
\right).
\]
\end{theorem}

\begin{proof}
By Lemma \ref{lem:product-localization},
\[
S^I[x^{-1}]
\cong
\prod_{i\in I}S[x_i^{-1}].
\]
If $i\in U(x)$, then $x_i$ is already a unit of $S$, so $S[x_i^{-1}]\cong S$.  If $i\in Z(x)$, Corollary \ref{cor:two-zero-localizations} gives $S[e^{-1}]\cong\B$.  If $i\in N(x)$, the same corollary gives $S[\tau^{-1}]\cong\zero$.  Since a finite product with a one-element factor is canonically isomorphic to the product of the remaining factors, the stated formula follows.  The explicit fraction formula is the coordinatewise description of these three identifications.
\end{proof}

\begin{example}
For $x=(\tau,e)\in S^2$,
\[
S^2[x^{-1}]\cong\B,
\]
and
\[
\frac{(a,b)}{(\tau,e)}=\chi_K(b).
\]
The first coordinate is erased; the second records only support.
\end{example}

\begin{example}
For $x=(h,e)\in S^2$ with $h\in K^\times$,
\[
S^2[x^{-1}]\cong S\times\B,
\]
and
\[
\frac{(a,b)}{(h,e)}=\bigl(ah^{-1},\chi_K(b)\bigr).
\]
This is ordinary division in one coordinate and Boolean division by supported zero in the other.
\end{example}

We next formalize the effect of adjoining an equality only on selected support coordinates.

For $i\in I$ and $s\in S$, write $r_i(s)\in S^I$ for the vector with entry $s$ at $i$ and $\tau$ elsewhere.  Fix a subset $A\subseteq I$ and units $h_i\in K^\times$ for $i\in A$.  Define
\[
\eps_A:=\sum_{i\in A}r_i(e),
\qquad
h_A:=\sum_{i\in A}r_i(h_i).
\]

\begin{theorem}[Scoped identity quotient]\label{thm:scoped-identity}
Let $\sim_A$ be the smallest semiring congruence on $S^I$ satisfying
\[
\eps_A\sim_A h_A.
\]
Then
\[
S^I/{\sim_A}
\cong
\B^A\times S^{I\setminus A}.
\]
The quotient map is
\[
q_A(x)
=
\left(
(\chi_K(x_i))_{i\in A},
(x_j)_{j\notin A}
\right).
\]
\end{theorem}

\begin{proof}
The map $q_A$ is a semiring homomorphism and satisfies
\[
q_A(\eps_A)=q_A(h_A),
\]
so it factors through $S^I/{\sim_A}$.

Let $f:S^I\to T$ be any semiring homomorphism satisfying $f(\eps_A)=f(h_A)$.  For $i\in A$, multiply the relation by $r_i(1)$ to obtain
\[
f(r_i(e))=f(r_i(h_i)).
\]
Multiplying again by $r_i(h_i^{-1})$ gives
\[
f(r_i(e))=f(r_i(1)).
\]
For every supported $s\in K$,
\[
r_i(s)=r_i(s)r_i(1),
\]
and therefore
\[
f(r_i(s))
=f(r_i(s))f(r_i(1))
=f(r_i(s))f(r_i(e))
=f(r_i(se))
=f(r_i(e)).
\]
Thus on coordinate $i\in A$, the value of $f$ depends only on whether the coordinate is $\tau$ or supported.  On coordinates outside $A$, no relation has been imposed.

Every vector $x\in S^I$ decomposes as
\[
x=\sum_{i\in I}r_i(x_i).
\]
Hence $f(x)$ depends only on $q_A(x)$, which gives a unique factorization of $f$ through $q_A$.  This is the universal property of the quotient.
\end{proof}

\begin{corollary}[Modal equality without arithmetic equality]\label{cor:modal-equality}
For every $h\in K^\times$,
\[
\chi_K(e)=\chi_K(h)=1.
\]
Thus $e$ and $h$ are equal after support reflection.  Their strict arithmetic identification gives the Boolean quotient by Theorem \ref{thm:strict-identification-collapse}; their identification on selected coordinates gives the partial Booleanization of Theorem \ref{thm:scoped-identity}.
\end{corollary}

\section{The split spectrum and the \texorpdfstring{$\mathbb F_1$}{F1} support skeleton}

\begin{theorem}[Ideals and prime ideals]\label{thm:prime-classification}
Let $R$ be a nonzero commutative ring and $S=G(R)$.
\begin{enumerate}[label=\textup{(\alph*)}]
\item Every ideal of $S$ is either $\{\tau\}$ or
\[
I_J:=J\cup\{\tau\}
\]
for a unique ideal $J\triangleleft R$.
\item The prime ideals are
\[
P_\tau:=\{\tau\}
\]
and
\[
P_{\mathfrak p}:=\mathfrak p\cup\{\tau\},
\qquad
\mathfrak p\in\Spec R.
\]
\end{enumerate}
\end{theorem}

\begin{proof}
Every ideal contains $\tau$.  If $I\cap R=\varnothing$, then $I=\{\tau\}$.  Otherwise $J:=I\cap R$ is closed under addition, multiplication by arbitrary elements of $R$, and additive inverses because $-a=(-1)a$.  Hence $J\triangleleft R$ and $I=J\cup\{\tau\}$.  Conversely every such set is an ideal.

The ideal $P_\tau$ is prime because a product in $S$ equals $\tau$ only when one factor equals $\tau$.  For $I_J$, the quotient by the corresponding ring-reflection congruence has supported part $R/J$; equivalently, for supported $a,b$,
\[
ab\in I_J
\iff
ab\in J.
\]
Thus $I_J$ is prime exactly when $J$ is prime.
\end{proof}

\begin{theorem}[Generic-point suspension]\label{thm:generic-suspension}
The map
\[
i:\Spec R\longrightarrow\Spec G(R),
\qquad
\mathfrak p\longmapsto P_{\mathfrak p},
\]
is a homeomorphism onto the closed subspace
\[
V(e)=\{P_{\mathfrak p}:\mathfrak p\in\Spec R\}.
\]
Furthermore,
\[
D(e)=\{P_\tau\},
\]
and $P_\tau$ is open, dense, and contained in every prime ideal.  Hence $\Spec G(R)$ is obtained from $\Spec R$ by adjoining a new open dense generic point below every classical prime in the specialization order.
\end{theorem}

\begin{proof}
Since $e=0_R$ belongs to every ring prime $\mathfrak p$ but does not belong to $P_\tau$, the displayed descriptions of $V(e)$ and $D(e)$ follow from Theorem \ref{thm:prime-classification}.  For an ideal $J\triangleleft R$,
\[
i(V_R(J))
=
V_S(I_J)\cap V_S(e),
\]
so $i$ identifies the Zariski closed sets and is a homeomorphism onto $V(e)$.

The containment $P_\tau\subseteq P$ holds for every prime $P$ because every ideal contains the semiring zero $\tau$.  Therefore the closure of $\{P_\tau\}$ is all of $\Spec S$.  Since $D(e)=\{P_\tau\}$, this point is also open.
\end{proof}

\begin{corollary}[Local stalks]\label{cor:split-stalks}
For $\mathfrak p\in\Spec R$,
\[
G(R)_{P_\tau}\cong\B,
\qquad
G(R)_{P_{\mathfrak p}}\cong G(R_{\mathfrak p}).
\]
\end{corollary}

\begin{proof}
The complement of $P_\tau$ is the whole supported part $R$, which contains $e$, so Corollary \ref{cor:fraction-comparison} gives $\B$.  The complement of $P_{\mathfrak p}$ is $R\setminus\mathfrak p$, which does not contain $0$; Theorem \ref{thm:ring-localization} gives $G(R_{\mathfrak p})$.
\end{proof}

\begin{example}[The shifted generic point of $\Spec\Z$]
For $S=G(\Z)$,
\[
P_\tau=\{\tau\}
\subsetneq
P_0=\{\tau,0\}
\subsetneq
P_p=p\Z\cup\{\tau\}
\]
for every rational prime $p$.  Thus the classical zero ideal $(0)$ becomes the nonzero prime $P_0$, while $P_\tau$ is the new zero ideal and new generic point.
\end{example}

\subsection{The exact blueprint presentation}

Let $M_S$ be the multiplicative monoid of $S$, including its absorbing element $\tau$.  Let $\N_0[M_S]$ denote the contracted monoid semiring in which the basis element $[\tau]$ is identified with the additive zero.

\begin{theorem}[Multiplicative skeleton plus preaddition]\label{thm:blueprint-presentation}
The canonical map
\[
\N_0[M_S]\longrightarrow S,
\qquad
[x]\longmapsto x,
\]
induces an isomorphism
\[
\N_0[M_S]\Big/\Big\langle [a]+[b]\sim[a+b]
\;:\;a,b\in R\Big\rangle
\xrightarrow{\sim}
G(R).
\]
Consequently the associated blueprint consists of the multiplicative monoid $M_S$ together with precisely the additive relations inherited from $R$ and the contraction of $\tau$ to additive zero.
\end{theorem}

\begin{proof}
The canonical map is surjective.  The imposed relations hold in $G(R)$, so the map factors through the quotient.

Every element of the quotient is represented by a finite sum of monoid generators.  Terms $[\tau]$ vanish in the contracted monoid semiring.  If at least one supported term remains, repeated use of
\[
[a]+[b]\sim[a+b]
\]
reduces the sum to a single generator $[r]$ with $r\in R$.  If no supported term remains, the sum is zero, represented by $[\tau]$.  Thus every congruence class has a representative in $M_S$.  Two such representatives have the same image in $G(R)$ only when they are equal, because the reduction relations reproduce the actual ring sum and never identify $\tau$ with a supported element.  Hence the induced map is bijective.
\end{proof}

\begin{remark}[The role of $\mathbb F_1$]
The monoid $M_S$ is the multiplicative $\mathbb F_1$ shadow in the sense of monoidal geometry, while Theorem \ref{thm:blueprint-presentation} gives the exact blueprint that restores additive cancellation data; see \cite{Deitmar,Lorscheid}.  The Boolean quotient $\chi_R:S\to\B$ is the pure support completion.  The projective geometry below makes this skeleton visible as a finite Boolean face poset whose fibres carry ordinary projective geometry.
\end{remark}

\section{Unimodular split projective geometry}

For a semiring $A$ and $n\ge0$, a vector
\[
x=(x_0,\dots,x_n)\in A^{n+1}
\]
is called \emph{unimodular} if the ideal generated by its coordinates is the unit ideal.  Equivalently, there exist $a_0,\dots,a_n\in A$ such that
\[
a_0x_0+\cdots+a_nx_n=1.
\]

\begin{definition}[Unimodular projective point set]\label{def:unimodular-projective}
Define
\[
\Pum^n(A)
:=
\{x\in A^{n+1}:x\text{ is unimodular}\}/A^\times,
\]
where $A^\times$ acts by simultaneous scalar multiplication.
\end{definition}

This is the usual set of projective points over a field.  Over a general semiring it is deliberately denoted by $\Pum$ rather than identified without qualification with a scheme-theoretic $\Proj$.

Let $K$ be a field and set $S=G(K)$.  For $x\in S^{n+1}$ define
\[
\supp(x):=\{i:x_i\in K\}.
\]
This records supported coordinates, including coordinates equal to $e=0_K$.

\begin{lemma}[Unimodularity criterion]\label{lem:unimodularity-criterion}
A vector $x\in S^{n+1}$ is unimodular if and only if at least one coordinate belongs to $K^\times$.
\end{lemma}

\begin{proof}
If $x_i\in K^\times$, then $x_i^{-1}x_i=1$, so the coordinates generate the unit ideal.

Conversely, suppose every coordinate belongs to $\{\tau,e\}$.  Multiplication of such a coordinate by any scalar remains in $\{\tau,e\}$, and a finite sum of elements of $\{\tau,e\}$ again lies in $\{\tau,e\}$.  It therefore cannot equal $1$.  Hence the vector is not unimodular.
\end{proof}

The support character induces a map
\[
\Supp:\Pum^n(S)\longrightarrow\Pum^n(\B),
\qquad
[x_0:\cdots:x_n]
\longmapsto
[\chi_K(x_0):\cdots:\chi_K(x_n)].
\]
Since $\B^\times=\{1\}$,
\[
\Pum^n(\B)=\B^{n+1}\setminus\{0\}
\]
is canonically the set of nonempty subsets of $[n]_0=\{0,\dots,n\}$.

\begin{theorem}[Boolean-face decomposition]\label{thm:projective-face-decomposition}
For every field $K$ and $n\ge0$, there is a canonical decomposition
\[
\boxed{
\Pum^n(G(K))
\cong
\bigsqcup_{\varnothing\neq A\subseteq[n]_0}
\mathbf P^{|A|-1}(K).}
\]
Under this identification, the fibre of $\Supp$ over $A$ is $\mathbf P^{|A|-1}(K)$.
\end{theorem}

\begin{proof}
Fix a nonempty subset $A\subseteq[n]_0$.  A vector of support exactly $A$ has the form
\[
x_i\in K\quad(i\in A),
\qquad
x_i=\tau\quad(i\notin A).
\]
By Lemma \ref{lem:unimodularity-criterion}, it is unimodular exactly when the amplitude vector $(x_i)_{i\in A}\in K^A$ is nonzero.  The unit group of $S$ is $K^\times$, and simultaneous multiplication by a supported unit acts on the amplitude vector by the usual scalar action.  Hence the set of projective classes of support $A$ is
\[
(K^A\setminus\{0\})/K^\times
=\mathbf P^{|A|-1}(K).
\]
Distinct support masks cannot be related by multiplication by a supported unit, so the union is disjoint.  The description of the support fibre is immediate.
\end{proof}

\begin{corollary}[Amplitude and support projections]\label{cor:projective-projections}
There are canonical maps
\[
\Pum^n(G(K))\xrightarrow{\;p_*\;}\mathbf P^n(K)
\]
and
\[
\Pum^n(G(K))\xrightarrow{\;\Supp\;}\Pum^n(\B).
\]
The first forgets whether a zero coordinate was absent or supported; the second forgets all amplitudes and retains only the support face.
\end{corollary}

\begin{proof}
A unimodular split vector has at least one nonzero amplitude coordinate, so its coordinatewise ring reflection is a nonzero vector in $K^{n+1}$.  This defines $p_*$.  The support map was defined above.  Both are invariant under multiplication by $K^\times$.
\end{proof}

\subsection{The split projective line}

For $n=1$, Theorem \ref{thm:projective-face-decomposition} gives
\[
\Pum^1(G(K))
\cong
\mathbf P^1(K)\sqcup\{0_\tau,\infty_\tau\},
\]
where
\[
0_\tau:=[\tau:1],
\qquad
\infty_\tau:=[1:\tau].
\]
Inside the full-support stratum $\mathbf P^1(K)$ lie the supported zero and supported infinity
\[
0_e:=[e:1],
\qquad
\infty_e:=[1:e].
\]

\begin{proposition}[Four zero-pole states]\label{prop:four-zero-pole-states}
The ring-reflection map satisfies
\[
p_*(0_\tau)=p_*(0_e)=0,
\qquad
p_*(\infty_\tau)=p_*(\infty_e)=\infty.
\]
The support map satisfies
\[
\Supp(0_\tau)=[0:1],
\qquad
\Supp(\infty_\tau)=[1:0],
\]
while every point of the full-support stratum, including $0_e$ and $\infty_e$, maps to
\[
[1:1].
\]
\end{proposition}

\begin{proof}
This follows directly from
\[
p_K(\tau)=p_K(e)=0,
\qquad
\chi_K(\tau)=0,
\quad
\chi_K(e)=1.
\]
\end{proof}

Thus the Boolean projective line
\[
\Pum^1(\B)=\{[0:1],[1:1],[1:0]\}
\]
is a three-point incidence skeleton: absent zero, full-support torus, and absent infinity.

\section{Finite-field counts and split projective symmetry}

\subsection{Point counts and the support-projective zeta function}

\begin{theorem}[Finite-field point count]\label{thm:finite-field-point-count}
For every prime power $q$ and every $n\ge0$,
\[
\#\Pum^n(G(\F_q))
=
\sum_{r=1}^{n+1}
\binom{n+1}{r}\frac{q^r-1}{q-1}
=
\frac{(q+1)^{n+1}-2^{n+1}}{q-1}.
\]
\end{theorem}

\begin{proof}
There are $\binom{n+1}{r}$ support faces of cardinality $r$, and the fibre over each such face is $\mathbf P^{r-1}(\F_q)$, which has $(q^r-1)/(q-1)$ points.  This gives the first formula.  For the second,
\[
\sum_{r=1}^{n+1}\binom{n+1}{r}(q^r-1)
=
\bigl((q+1)^{n+1}-1\bigr)-\bigl(2^{n+1}-1\bigr),
\]
which equals $(q+1)^{n+1}-2^{n+1}$.
\end{proof}

\begin{definition}[Support-projective zeta function]\label{def:support-zeta}
Define
\[
Z_{\mathrm{sp}}(\Pum^n/\F_q,T)
:=
\exp\left(
\sum_{m\ge1}
\#\Pum^n(G(\F_{q^m}))\frac{T^m}{m}
\right).
\]
This is a point-counting zeta function for the split projective functor; no assertion is made here that $\Pum^n(G(\F_q))$ is the rational-point set of an ordinary scheme.
\end{definition}

\begin{theorem}[Closed zeta product]\label{thm:support-zeta-product}
Put
\[
c_{n,j}:=\sum_{r=j+1}^{n+1}\binom{n+1}{r}
\qquad(0\le j\le n).
\]
Then
\[
\#\Pum^n(G(\F_q))
=
\sum_{j=0}^{n}c_{n,j}q^j,
\]
and
\[
\boxed{
Z_{\mathrm{sp}}(\Pum^n/\F_q,T)
=
\prod_{j=0}^{n}(1-q^jT)^{-c_{n,j}}.}
\]
In particular,
\[
Z_{\mathrm{sp}}(\Pum^1/\F_q,T)
=
\frac{1}{(1-T)^3(1-qT)}.
\]
\end{theorem}

\begin{proof}
Since
\[
\frac{q^r-1}{q-1}=1+q+\cdots+q^{r-1},
\]
the coefficient of $q^j$ in the point count of Theorem \ref{thm:finite-field-point-count} is the sum of $\binom{n+1}{r}$ over $r\ge j+1$, which is $c_{n,j}$.

Substituting $q^m$ for $q$ gives
\[
\#\Pum^n(G(\F_{q^m}))
=
\sum_{j=0}^{n}c_{n,j}q^{mj}.
\]
Therefore
\[
\begin{aligned}
\log Z_{\mathrm{sp}}
&=
\sum_{m\ge1}\sum_{j=0}^{n}c_{n,j}q^{mj}\frac{T^m}{m}\\
&=
\sum_{j=0}^{n}c_{n,j}
\sum_{m\ge1}\frac{(q^jT)^m}{m}\\
&=
-\sum_{j=0}^{n}c_{n,j}\log(1-q^jT),
\end{aligned}
\]
which exponentiates to the product formula.  For $n=1$,
\[
c_{1,0}=\binom21+\binom22=3,
\qquad
c_{1,1}=\binom22=1.
\]
\end{proof}

\begin{remark}
The ordinary projective-line zeta function is
\[
Z(\mathbf P^1/\F_q,T)=\frac{1}{(1-T)(1-qT)}.
\]
The split support completion contributes two additional degree-zero factors, corresponding to $0_\tau$ and $\infty_\tau$.
\end{remark}

\subsection{Invertible matrices}

\begin{lemma}[Boolean invertible matrices]\label{lem:boolean-gl}
Every invertible $m\times m$ matrix over $\B$ is a permutation matrix.  Hence
\[
\GL_m(\B)\cong S_m.
\]
\end{lemma}

\begin{proof}
The free $\B$-semimodule $\B^m$ is the Boolean lattice of subsets of $\{1,\dots,m\}$, with addition given by union.  A $\B$-linear automorphism preserves joins and zero, hence is an automorphism of this finite Boolean lattice.  The atoms are exactly the standard basis vectors.  Every lattice automorphism permutes the atoms, so the corresponding matrix is a permutation matrix.  Conversely every permutation matrix is invertible.
\end{proof}

\begin{theorem}[Monomial general linear group]\label{thm:monomial-gl}
Let $K$ be a field and $S=G(K)$.  Then every matrix in $\GL_m(S)$ is monomial with one entry from $K^\times$ in each row and column and $\tau$ in every other position.  Consequently,
\[
\boxed{
\GL_m(G(K))\cong(K^\times)^m\rtimes S_m.}
\]
\end{theorem}

\begin{proof}
Let $U\in\GL_m(S)$ and let $V$ be its inverse.  Applying the support character entrywise gives
\[
\chi(U)\chi(V)=I_m,
\qquad
\chi(V)\chi(U)=I_m
\]
over $\B$.  By Lemma \ref{lem:boolean-gl}, $\chi(U)$ is a permutation matrix.  Therefore $U$ has exactly one supported entry in each row and column and has $\tau$ elsewhere.

Applying the ring reflection entrywise gives
\[
p(U)p(V)=I_m
\]
over $K$.  Thus $p(U)$ is an invertible monomial matrix, so each of its monomial entries is nonzero.  The corresponding supported entries of $U$ therefore lie in $K^\times$.

Conversely, a monomial matrix with entries in $K^\times$ has the evident monomial inverse.  The semidirect-product description is the standard decomposition into diagonal and permutation matrices.
\end{proof}

\begin{corollary}[Split projective linear group]\label{cor:split-pgl}
The group of projective transformations induced by invertible split matrices is
\[
\operatorname{PGL}^{\mathrm{lin}}_m(G(K))
\cong
\bigl((K^\times)^m/K^\times\bigr)\rtimes S_m,
\]
where the denominator is the diagonal scalar subgroup.  For $m=2$,
\[
\operatorname{PGL}^{\mathrm{lin}}_2(G(K))
\cong
K^\times\rtimes C_2.
\]
On the full-support affine coordinate $z$, the transformations are exactly
\[
z\longmapsto az
\qquad\text{and}\qquad
z\longmapsto\frac{a}{z},
\qquad a\in K^\times.
\]
\end{corollary}

\begin{proof}
Projectivization quotients $\GL_m(G(K))$ by the central diagonal copy of $K^\times$.  For $m=2$, a diagonal monomial matrix acts by scaling and an antidiagonal monomial matrix acts by scaling followed by inversion.
\end{proof}

\begin{corollary}[Obstruction to global translation]\label{cor:no-global-translation}
For $b\neq0$, the classical translation
\[
z\longmapsto z+b
\]
does not arise from a global split-projective linear automorphism of $\Pum^1(G(K))$.  More generally, a classical M\"obius transformation extends linearly to the mixed-support projective line only when its matrix is monomial after replacing structural zeros by $\tau$.
\end{corollary}

\begin{proof}
A matrix for translation is
\[
\begin{pmatrix}1&b\\0&1\end{pmatrix}.
\]
Its support pattern is not a permutation matrix, regardless of whether the classical zero entry is represented by $e$ or by $\tau$.  Theorem \ref{thm:monomial-gl} therefore excludes invertibility over $G(K)$.  The same argument applies to a general M\"obius matrix.
\end{proof}

\begin{remark}[Support-homogeneous affine maps]
Corollary \ref{cor:no-global-translation} does not say that affine maps disappear.  It says that they cannot be represented by the naive homogeneous coordinate $(x,1)$ on all support fibres.  The support-idempotent homogenization of \cite{Homogenization} replaces $(x,1)$ by $(x,h)$, where $h$ is the support idempotent, and represents
\[
x\longmapsto ax+b
\]
by
\[
(x,h)\longmapsto(ax+bh,h).
\]
This is the correct fibrewise extension across the Boolean support diagram.
\end{remark}

\section{Residue monads and Boolean-cubical infinitesimal geometry}

The support residue $\B$ is not itself Leibniz's infinitesimal monad.  A monad requires a residue map with a nontrivial kernel of infinitesimal amplitudes.  Split zero supplies the support geometry in which such kernels are organized.

\begin{definition}[Split residue datum]\label{def:split-residue-datum}
A \emph{split residue datum} is a surjective ring homomorphism
\[
q:A\twoheadrightarrow k
\]
with kernel
\[
\mathfrak m:=\ker q.
\]
It induces a semiring homomorphism
\[
G(q):G(A)\longrightarrow G(k)
\]
by
\[
G(q)(\tau)=\tau,
\qquad
G(q)(a)=q(a)
\quad(a\in A).
\]
\end{definition}

\begin{definition}[Supported monad]
For a supported element $a\in k\subset G(k)$ and a chosen lift $\widetilde a\in A$, define
\[
\mu_q(a):=G(q)^{-1}(a).
\]
For the absent state define
\[
\mu_q(\tau):=G(q)^{-1}(\tau).
\]
\end{definition}

\begin{theorem}[Monad fibre formula]\label{thm:monad-fibre}
Let $q:A\twoheadrightarrow k$ be a split residue datum.  Then
\[
\mu_q(a)=\widetilde a+\mathfrak m
\qquad(a\in k),
\]
and
\[
\mu_q(\tau)=\{\tau\}.
\]
In particular,
\[
\mu_q(e_k)=\mathfrak m,
\]
where $e_k=0_k$ is the supported zero in $G(k)$.
\end{theorem}

\begin{proof}
A supported element $x\in A\subset G(A)$ maps to $a$ exactly when $q(x)=a$, which is equivalent to $x\in\widetilde a+\mathfrak m$.  No supported element maps to $\tau$, because $G(q)$ preserves support.  Hence the only preimage of $\tau$ is $\tau$ itself.
\end{proof}

\begin{corollary}[Centre and halo]\label{cor:centre-and-halo}
The supported zero $e_k$ is the centre of the zero monad, the ideal $\mathfrak m$ is its infinitesimal halo, and $\tau$ is the absent state outside that halo.
\end{corollary}

\begin{proof}
The fibre over $e_k$ is $\mathfrak m$ by Theorem \ref{thm:monad-fibre}.  Its distinguished zero element is $0_A=e_A$, which maps to $e_k$.  The absent state has the singleton fibre $\{\tau\}$.
\end{proof}

The higher-dimensional construction is cubical.  Let $I$ be finite and let
\[
a=(a_i)_{i\in I}\in G(k)^I.
\]
Write
\[
J:=\supp(a)=\{i:a_i\in k\}.
\]
For $i\in J$, choose a lift $\widetilde a_i\in A$.

\begin{definition}[Support star and cubical monad]\label{def:cubical-monad}
The \emph{support star} of $a$ is
\[
\operatorname{Star}(a)
:=
\left\{
 b\in G(k)^I:
 b_i=a_i\ (i\in J),
\quad
b_i\in\{\tau,e_k\}\ (i\notin J)
\right\}.
\]
The \emph{upward cubical monad} of $a$ is
\[
\mu_q^\square(a)
:=
(G(q)^I)^{-1}(\operatorname{Star}(a)).
\]
\end{definition}

\begin{theorem}[Boolean-cubical monad decomposition]\label{thm:cubical-monad}
There is a canonical disjoint decomposition
\[
\boxed{
\mu_q^\square(a)
=
\bigsqcup_{L\subseteq I\setminus J}
\left(
\prod_{i\in J}(\widetilde a_i+\mathfrak m)
\right)
\times
\mathfrak m^L
\times
\{\tau\}^{I\setminus(J\cup L)}.}
\]
The subset $L$ records exactly the formerly absent coordinates that have been activated with residue zero.
\end{theorem}

\begin{proof}
Every $b\in\operatorname{Star}(a)$ is determined by the subset
\[
L:=\{i\in I\setminus J:b_i=e_k\}.
\]
For $i\in J$, Theorem \ref{thm:monad-fibre} gives the fibre $\widetilde a_i+\mathfrak m$.  For $i\in L$, the fibre of $e_k$ is $\mathfrak m$.  For the remaining coordinates, the fibre of $\tau$ is $\{\tau\}$.  Different subsets $L$ produce disjoint support patterns, so the union is disjoint.
\end{proof}

\begin{example}[The local path from absence to infinitesimal amplitude]
Take one coordinate and the absent base point $a=\tau$.  Then
\[
\mu_q^\square(\tau)
=
\{\tau\}\sqcup\mathfrak m.
\]
Inside the second component lie
\[
e_A=0_A
\]
and all nonzero infinitesimals $h\in\mathfrak m\setminus\{0\}$.  Thus the exact local sequence is
\[
\tau
\longrightarrow
 e_A
\longrightarrow
 h\in\mathfrak m\setminus\{0\},
\]
meaning absence, supported exact zero, and supported infinitesimal amplitude.
\end{example}

\subsection{First-order cubical tangent data}

Suppose now that $k$ is a field.  The ordinary first-order tangent space of the residue datum is the $k$-vector space
\[
T_q:=\mathfrak m/\mathfrak m^2.
\]
The cubical monad carries one copy of this vector space for every active coordinate.

\begin{theorem}[Boolean cube of first-order tangent fibres]\label{thm:first-order-cubical}
Let $q:A\twoheadrightarrow k$ be a surjection onto a field, let $a\in G(k)^I$ have support $J$, and fix a stratum indexed by $L\subseteq I\setminus J$ in Theorem \ref{thm:cubical-monad}.  Modulo coordinatewise congruence by $\mathfrak m^2$, that stratum is an affine torsor under
\[
T_q^{J\cup L}.
\]
As $L$ varies, the first-order split neighbourhood is a Boolean cube of affine tangent spaces indexed by $\mathcal P(I\setminus J)$.
\end{theorem}

\begin{proof}
For each $i\in J$, the coset $\widetilde a_i+\mathfrak m$ modulo $\mathfrak m^2$ is an affine torsor under $\mathfrak m/\mathfrak m^2$.  For each newly activated coordinate $i\in L$, the fibre $\mathfrak m$ modulo $\mathfrak m^2$ is the vector space $T_q$ itself.  Coordinates outside $J\cup L$ are fixed at $\tau$.  Taking products gives an affine torsor under $T_q^{J\cup L}$.  The indexing subsets $L$ form the Boolean cube of possible support activations.
\end{proof}

\begin{remark}
Theorem \ref{thm:first-order-cubical} separates two independent infinitesimal directions:
\[
\text{Boolean activation between support faces}
\]
and
\[
\text{linear tangent displacement inside an active face}.
\]
The element $e$ controls the first, while $\mathfrak m/\mathfrak m^2$ controls the second.
\end{remark}

\subsection{Formal completion}

Let $A$ be a ring and $I\triangleleft A$ an ideal.  The quotient maps $A/I^{r+1}\to A/I^r$ induce an inverse system of split-zero semirings.

\begin{theorem}[Split zero commutes with adic completion]\label{thm:completion-commutes}
There is a canonical semiring isomorphism
\[
\varprojlim_{r\ge1}G(A/I^r)
\xrightarrow{\sim}
G\left(\varprojlim_{r\ge1}A/I^r\right).
\]
In particular, if $A$ is $I$-adically complete, then
\[
\varprojlim_{r\ge1}G(A/I^r)
\cong
G(A).
\]
\end{theorem}

\begin{proof}
An element of the inverse limit on the left is a coherent sequence $(x_r)_{r\ge1}$.  Every transition map sends $\tau$ to $\tau$ and sends supported elements to supported elements.  Hence either all $x_r$ equal $\tau$, or every $x_r$ is supported.  In the first case the sequence corresponds to the external zero of the right-hand side.  In the second case the sequence is exactly a coherent sequence in the ring quotients $A/I^r$, hence an element of $\varprojlim A/I^r$, regarded as supported.  This gives a bijection, and the operations are preserved coordinatewise.
\end{proof}

\begin{definition}[Split formal monad]
The split formal neighbourhood associated with $(A,I)$ is the pro-semiring
\[
\operatorname{Spf}_{\mathrm{sp}}(A,I)
:=
\{G(A/I^r)\}_{r\ge1}.
\]
Its supported sector is the ordinary formal neighbourhood; its support masks form the Boolean incidence skeleton.
\end{definition}

\section{Hyperreal and surcomplex realizations}

\subsection{Finite hyperreals}

Let ${}^{\ast}\R$ be a nonstandard elementary extension of $\R$.  Let
\[
\cO_{\mathrm{hyp}}
:=
\{x\in{}^{\ast}\R:|x|<n\text{ for some }n\in\N\}
\]
be the ring of finite hyperreals and let
\[
\mathfrak m_{\mathrm{hyp}}
:=
\{x\in{}^{\ast}\R:|x|<1/n\text{ for every }n\ge1\}
\]
be its ideal of infinitesimals.  The standard-part map is a surjective ring homomorphism
\[
\st:\cO_{\mathrm{hyp}}\twoheadrightarrow\R
\]
with kernel $\mathfrak m_{\mathrm{hyp}}$; see \cite{Robinson,Goldblatt}.

\begin{corollary}[Split hyperreal monads]\label{cor:hyperreal-monads}
For every $a\in\R$,
\[
\mu_{\mathrm{hyp}}(a)
=
a+\mathfrak m_{\mathrm{hyp}}.
\]
For a mixed-support point $a\in G(\R)^I$, the upward monad is the Boolean cube of ordinary hyperreal monads described by Theorem \ref{thm:cubical-monad}.
\end{corollary}

\begin{proof}
Apply Theorems \ref{thm:monad-fibre} and \ref{thm:cubical-monad} to the standard-part map.
\end{proof}

\subsection{Surreal and surcomplex normal forms}

We work in von Neumann--Bernays--G\"odel set theory so that the proper class $\No$ may be treated as a class-sized ordered field.  Every support occurring in a Conway normal form is a set.  We use Conway's normal-form theorem and the standard monomial map
\[
x\longmapsto\omega^x
\]
from \cite{Conway}.  The uploaded expository notes \cite{Simons} give a concise account of birthdays, ordinal embeddings, and the first infinitesimal $\omega^{-1}$.

Set
\[
K:=\No[i]
=
\{x+iy:x,y\in\No,\ i^2=-1\}.
\]
Conway proved that $K$ is algebraically closed \cite{Conway}.

\begin{theorem}[Surcomplex Conway normal form]\label{thm:surcomplex-normal-form}
Every nonzero $z\in K$ has a unique expansion
\[
z=
\sum_{\rho<\lambda}
 c_\rho\omega^{\gamma_\rho},
\]
where
\[
c_\rho\in\C^\times,
\qquad
\gamma_0>\gamma_1>\cdots,
\]
and the support $\{\gamma_\rho:\rho<\lambda\}$ is reverse well ordered.  Equivalently, $K$ is the class-sized Hahn field with coefficient field $\C$ and value group $\No$, with set-sized reverse-well-ordered support.
\end{theorem}

\begin{proof}
Write $z=x+iy$ with $x,y\in\No$.  By Conway's normal-form theorem,
\[
x=\sum_{\gamma\in A}a_\gamma\omega^\gamma,
\qquad
 y=\sum_{\gamma\in B}b_\gamma\omega^\gamma,
\]
where $A$ and $B$ are reverse well ordered and the coefficients are real.  The finite union $A\cup B$ is reverse well ordered.  Extending missing coefficients by zero gives
\[
z=
\sum_{\gamma\in A\cup B}
(a_\gamma+ib_\gamma)\omega^\gamma.
\]
Delete zero coefficients.  Uniqueness follows by comparing real and imaginary parts and using uniqueness of the real Conway normal form.
\end{proof}

For nonzero $z$, let $\ell(z)$ be its leading exponent and define the natural valuation
\[
v(z):=-\ell(z),
\qquad
v(0):=\infty.
\]

\begin{proposition}[Valuation laws]\label{prop:surcomplex-valuation}
For $z,w\in K$,
\[
v(zw)=v(z)+v(w),
\]
and
\[
v(z+w)\ge\min\{v(z),v(w)\},
\]
with equality unless the leading terms cancel.  If $z=x+iy$, then
\[
v(z)=\min\{v(x),v(y)\}.
\]
\end{proposition}

\begin{proof}
If
\[
z=c\omega^\alpha+\text{lower terms},
\qquad
w=d\omega^\beta+\text{lower terms},
\]
then
\[
zw=cd\,\omega^{\alpha+\beta}+\text{lower terms},
\]
and $cd\neq0$.  Thus $\ell(zw)=\alpha+\beta$, proving multiplicativity of $v$.  The addition inequality is the standard leading-term argument.

For $z=x+iy$, let the leading real exponents of $x$ and $y$ be $\alpha$ and $\beta$.  If $\alpha\neq\beta$, the larger exponent survives in $z$.  If $\alpha=\beta$, the leading complex coefficient is $a+ib$ with $a,b\in\R$ not both zero, so there is no cancellation between real and imaginary components.  Hence $\ell(z)=\max\{\alpha,\beta\}$, which is equivalent to the displayed valuation formula.
\end{proof}

Define the finite valuation ring and its maximal ideal by
\[
\cO_K:=\{z\in K:v(z)\ge0\},
\qquad
\mathfrak m_K:=\{z\in K:v(z)>0\}.
\]

\begin{theorem}[Surcomplex standard part]\label{thm:surcomplex-standard-part}
The coefficient of $\omega^0$ defines a surjective ring homomorphism
\[
\st:\cO_K\twoheadrightarrow\C
\]
with kernel $\mathfrak m_K$.  Hence
\[
\cO_K/\mathfrak m_K\cong\C.
\]
\end{theorem}

\begin{proof}
An element of $\cO_K$ has no positive exponent in its normal form.  Define $\st(z)$ to be the coefficient of $\omega^0$, taking that coefficient to be zero if the exponent $0$ is absent.  Addition preserves constant coefficients.  In a product of two finite series, an exponent sum can equal zero only when both exponents are zero, because all exponents are nonpositive.  Therefore the constant coefficient of a product is the product of the constant coefficients, so $\st$ is a ring homomorphism.  It is surjective because every $c\in\C$ is the constant series $c\omega^0$.  Its kernel consists exactly of series whose leading exponent is negative, equivalently those of positive valuation.
\end{proof}

\begin{corollary}[The surcomplex Leibniz monad]\label{cor:surcomplex-monad}
For every $c\in\C$,
\[
\mu_{\mathrm{sur}}(c)
=
c+\mathfrak m_K.
\]
In particular, the zero monad is
\[
\mu_{\mathrm{sur}}(0)=\mathfrak m_K.
\]
After split-zero globalization, the upward mixed-support monad is the Boolean cube of these surcomplex fibres.
\end{corollary}

\begin{proof}
Apply Theorems \ref{thm:monad-fibre} and \ref{thm:cubical-monad} to the standard-part map of Theorem \ref{thm:surcomplex-standard-part}.
\end{proof}

\subsection{Ordinal scales and reciprocal infinitesimals}

Every ordinal embeds canonically in $\No$.  Inside the ordered field $\No$, its arithmetic is the natural or Hessenberg arithmetic induced by Conway normal forms, not the noncommutative ordinal addition and multiplication.

\begin{proposition}[Reciprocals of infinite ordinals]\label{prop:ordinal-reciprocals}
If $\alpha$ is an infinite ordinal, regarded as a positive surreal number, then
\[
\alpha^{-1}
\]
is a positive surreal infinitesimal.  The map
\[
\alpha\longmapsto\alpha^{-1}
\]
is order reversing on positive ordinals.  Since $\omega$ is the least infinite ordinal, $\omega^{-1}$ is the largest reciprocal of an infinite ordinal, although it is not the largest positive surreal infinitesimal.
\end{proposition}

\begin{proof}
For every positive integer $n$,
\[
\alpha>n.
\]
Inversion in an ordered field reverses inequalities among positive elements, so
\[
0<\alpha^{-1}<n^{-1}.
\]
This is the definition of a positive infinitesimal.  Order reversal is the same argument.  The final assertion follows from minimality of $\omega$ among infinite ordinals.  There is no largest positive infinitesimal because $2\epsilon$ is infinitesimal and larger than $\epsilon$ whenever $\epsilon>0$ is infinitesimal.
\end{proof}

\begin{theorem}[Reciprocal ordinals are coinitial]\label{thm:reciprocal-coinitial}
For every positive surreal infinitesimal $\epsilon$, there exists an infinite ordinal $\alpha$ such that
\[
0<\alpha^{-1}<\epsilon.
\]
\end{theorem}

\begin{proof}
The reciprocal $x:=\epsilon^{-1}$ is positive infinite.  The ordinals are cofinal in $\No$; equivalently, every surreal number is bounded above by an ordinal \cite[Chap.~5]{Conway}.  Choose an ordinal $\alpha>x$.  Then $\alpha$ is infinite and inversion gives
\[
0<\alpha^{-1}<x^{-1}=\epsilon.
\]
\end{proof}

\begin{lemma}[Geometric inversion in the Hahn field]\label{lem:hahn-geometric}
If $u\in K$ satisfies $v(u)>0$, then the family $\{(-u)^n:n\ge0\}$ is Hahn-summable and
\[
(1+u)^{-1}=\sum_{n=0}^{\infty}(-u)^n.
\]
\end{lemma}

\begin{proof}
Since $v(u^n)=nv(u)$ tends strictly upward, for every valuation bound only finitely many terms have valuation below that bound.  This is precisely the Hahn summability condition.  Multiplying partial sums gives
\[
(1+u)\sum_{n=0}^{N}(-u)^n
=1+(-1)^Nu^{N+1}.
\]
The remainder tends to arbitrarily high valuation, so the Hahn sum is an inverse of $1+u$.
\end{proof}

\begin{theorem}[Reciprocal from Cantor normal form]\label{thm:ordinal-inverse-expansion}
Let an infinite ordinal have Cantor normal form
\[
\alpha
=
n_0\omega^{\beta_0}
+n_1\omega^{\beta_1}
+\cdots
+n_m\omega^{\beta_m},
\]
where
\[
\beta_0>\beta_1>\cdots>\beta_m,
\qquad
n_j\in\N_{>0}.
\]
Set
\[
u:=
\sum_{j=1}^{m}
\frac{n_j}{n_0}\omega^{\beta_j-\beta_0}.
\]
Then $u$ is infinitesimal and
\[
\boxed{
\alpha^{-1}
=
\frac1{n_0}\omega^{-\beta_0}
\sum_{k=0}^{\infty}(-u)^k.}
\]
Its leading term is
\[
\frac1{n_0}\omega^{-\beta_0},
\]
so
\[
v(\alpha^{-1})=\beta_0.
\]
\end{theorem}

\begin{proof}
Inside $\No$, Cantor normal form is Conway normal form for ordinals.  Factor
\[
\alpha
=
n_0\omega^{\beta_0}(1+u).
\]
Each exponent $\beta_j-\beta_0$ is negative, so $v(u)>0$.  Lemma \ref{lem:hahn-geometric} gives the inverse of $1+u$, and multiplication by $(n_0\omega^{\beta_0})^{-1}$ gives the formula.  Since the geometric factor has leading term $1$, the leading term of $\alpha^{-1}$ is the displayed monomial.
\end{proof}

\begin{example}
For $\alpha=\omega+1$,
\[
\frac1{\omega+1}
=
\omega^{-1}
-\omega^{-2}
+\omega^{-3}
-\omega^{-4}
+\cdots.
\]
This is an exact surreal equality.
\end{example}

\subsection{Complex phase at every infinitesimal order}

For $\lambda\in\No$, define
\[
F_{\ge\lambda}:=\{z\in K:v(z)\ge\lambda\},
\qquad
F_{>\lambda}:=\{z\in K:v(z)>\lambda\}.
\]

\begin{theorem}[Complex graded layers]\label{thm:complex-graded-layers}
For every $\lambda\in\No$, coefficient extraction at the monomial $\omega^{-\lambda}$ induces a canonical $\C$-linear isomorphism
\[
F_{\ge\lambda}/F_{>\lambda}
\xrightarrow{\sim}
\C.
\]
Multiplication of initial terms gives bilinear maps
\[
\gr_\lambda(K)\times\gr_\mu(K)
\longrightarrow
\gr_{\lambda+\mu}(K)
\]
corresponding to multiplication in $\C$.
\end{theorem}

\begin{proof}
For $z\in F_{\ge\lambda}$, let $c_{-\lambda}(z)$ be the coefficient of $\omega^{-\lambda}$, taking it to be zero when that monomial is absent.  This is a $\C$-linear map to $\C$.  Its kernel consists exactly of elements with valuation strictly greater than $\lambda$, namely $F_{>\lambda}$.  It is surjective because $c\omega^{-\lambda}$ maps to $c$.  The multiplication statement follows from
\[
(c\omega^{-\lambda})(d\omega^{-\mu})
=cd\,\omega^{-(\lambda+\mu)}.
\]
\end{proof}

\begin{definition}[Support--order--phase initial datum]\label{def:support-order-phase}
For $x\in G(K)$ define
\[
\operatorname{in}_{\mathrm{sp}}(x)
=
\begin{cases}
(0,\bot,0),&x=\tau,\\
(1,\infty,0),&x=e,\\
(1,v(x),c),&x\in K^\times,
\end{cases}
\]
where $c\in\C^\times$ is the leading coefficient of $x$.  Here $\bot$ denotes absence and $\infty$ denotes supported exact zero.
\end{definition}

\begin{remark}
The three coordinates have independent meanings:
\[
\text{support}\in\B,
\qquad
\text{order}\in\No\cup\{\infty\},
\qquad
\text{phase}\in\C.
\]
The split-zero residue retains the first, the valuation retains the second, and the associated graded retains the third.
\end{remark}

\begin{corollary}[Mixed surreal division]\label{cor:mixed-surreal-division}
Let $\alpha>0$ be an ordinal and let $S=G(\No[i])$.  Then
\[
S^2[(\omega^\alpha,e)^{-1}]
\cong
S\times\B,
\]
and
\[
\frac{(a,b)}{(\omega^\alpha,e)}
=
\bigl(a\omega^{-\alpha},\chi(b)\bigr).
\]
Similarly,
\[
\frac{(a,b)}{(\omega^{-\alpha},e)}
=
\bigl(a\omega^{\alpha},\chi(b)\bigr).
\]
\end{corollary}

\begin{proof}
Every nonzero surreal or surcomplex number is a unit of $K$.  Apply Theorem \ref{thm:mixed-localization}.
\end{proof}

\begin{corollary}[Strictly identifying $e$ with a surreal infinitesimal]
For every nonzero $h\in\No[i]$, including $h=\omega^{-\alpha}$,
\[
G(\No[i])/(e=h)\cong\B.
\]
On a selected coordinate face $A$, the corresponding quotient is the partial Booleanization of Theorem \ref{thm:scoped-identity}.
\end{corollary}

\begin{proof}
Apply Theorems \ref{thm:strict-identification-collapse} and \ref{thm:scoped-identity}.
\end{proof}

\subsection{Differential separation of $e$ and $\omega^{-1}$}

Berarducci and Mantova construct a simplest surreal derivation
\[
\partial:\No\longrightarrow\No
\]
with constant field $\R$, strong additivity, compatibility with exponentiation, and
\[
\partial(\omega)=1;
\]
see \cite{BerarducciMantova}.  Extend it $\C$-linearly to $K=\No[i]$ by $\partial(i)=0$.

\begin{proposition}[Differential obstruction to identification]\label{prop:differential-obstruction}
One has
\[
\partial(\omega^{-1})=-\omega^{-2}\neq0,
\]
whereas every semiring derivation on $G(K)$ kills the supported zero $e$.  Hence no derivation-preserving identification can send $e$ to $\omega^{-1}$.
\end{proposition}

\begin{proof}
Differentiate $\omega\omega^{-1}=1$:
\[
0
=
\partial(1)
=
\partial(\omega)\omega^{-1}
+
\omega\partial(\omega^{-1})
=
\omega^{-1}
+
\omega\partial(\omega^{-1}).
\]
Multiplication by $\omega^{-1}$ gives
\[
\partial(\omega^{-1})=-\omega^{-2}.
\]
The statement for $e$ is Proposition \ref{prop:derivation-kills-e}.
\end{proof}

\section{Projective zero, projective infinity, and the surcomplex circle}

\subsection{Inversion and the two poles}

Let $K$ be a field and $S=G(K)$.  Consider the matrix
\[
J:=
\begin{pmatrix}
\tau&1\\
1&\tau
\end{pmatrix}.
\]

\begin{proposition}[Global split inversion]\label{prop:global-inversion}
The matrix $J$ belongs to $\GL_2(S)$ and satisfies $J^2=I$.  Its action on $\Pum^1(S)$ is
\[
[x:y]\longmapsto[y:x].
\]
In particular,
\[
0_\tau\longleftrightarrow\infty_\tau,
\qquad
0_e\longleftrightarrow\infty_e.
\]
\end{proposition}

\begin{proof}
A direct multiplication gives
\[
J^2
=
\begin{pmatrix}
\tau^2+1&\tau+\tau\\
\tau+\tau&1+\tau^2
\end{pmatrix}
=
\begin{pmatrix}
1&\tau\\
\tau&1
\end{pmatrix}
=I.
\]
The coordinate action and the stated exchanges are immediate.
\end{proof}

\begin{theorem}[Equivalence of zero and infinity monads]\label{thm:zero-infinity-monads}
Let $q:A\twoheadrightarrow k$ be a split residue datum and let $S_A=G(A)$.  In the affine chart $[x:1]$, the upward zero monad is
\[
\{[\tau:1]\}
\sqcup
\{[h:1]:h\in\mathfrak m\}.
\]
Projective inversion restricts to a bijection onto the upward infinity monad
\[
\{[1:\tau]\}
\sqcup
\{[1:h]:h\in\mathfrak m\}.
\]
It sends
\[
[\tau:1]\mapsto[1:\tau],
\qquad
[e:1]\mapsto[1:e],
\]
and, for $h\neq0$,
\[
[h:1]\mapsto[1:h]=[h^{-1}:1]
\]
whenever $A$ is contained in a field in which $h$ is invertible.
\end{theorem}

\begin{proof}
The description of the upward zero monad is Theorem \ref{thm:cubical-monad} in one coordinate, written in the affine projective chart.  Proposition \ref{prop:global-inversion} swaps the homogeneous coordinates and gives the displayed image.  If $h$ is nonzero in an ambient field, multiplication by the unit $h^{-1}$ identifies $[1:h]$ with $[h^{-1}:1]$.
\end{proof}

\begin{corollary}[No equation $x=\infty$ is required]\label{cor:no-x-infinity}
The supported projective infinity
\[
\infty_e=[1:e]
\]
and the absent-denominator infinity
\[
\infty_\tau=[1:\tau]
\]
are points of $\Pum^1(G(K))$.  Infinite field elements form the punctured valuative monad of $\infty_e$ under inversion; they are not obtained by adjoining an arithmetic element satisfying $x=\infty$.
\end{corollary}

\begin{proof}
The two points are distinct by their support masks.  The punctured-monad assertion is Theorem \ref{thm:zero-infinity-monads} applied to a valued field: nonzero infinitesimals invert to infinite elements.
\end{proof}

\subsection{The Cayley circle}

Let $F$ be a real closed field, let $F[i]$ be its algebraic closure, and define
\[
U_F:=\{z\in F[i]:z\overline z=1\}.
\]

\begin{theorem}[Algebraic Cayley circle]\label{thm:cayley-circle}
The Cayley transform
\[
C:\mathbf P^1(F)\longrightarrow U_F
\]
defined by
\[
C(x)=\frac{x-i}{x+i}
\qquad(x\in F),
\qquad
C(\infty)=1,
\]
is a bijection.  Its inverse is
\[
C^{-1}(z)=i\frac{1+z}{1-z}
\qquad(z\neq1),
\qquad
C^{-1}(1)=\infty.
\]
In particular,
\[
\mathbf P^1(\No)
\cong
\{z\in\No[i]:z\overline z=1\}.
\]
\end{theorem}

\begin{proof}
For $x\in F$,
\[
C(x)\overline{C(x)}
=
\frac{x-i}{x+i}\frac{x+i}{x-i}=1.
\]
Conversely, let $z\in U_F$ with $z\neq1$.  Since $\overline z=z^{-1}$,
\[
\overline{\left(i\frac{1+z}{1-z}\right)}
=
-i\frac{1+z^{-1}}{1-z^{-1}}
=
-i\frac{z+1}{z-1}
=
i\frac{1+z}{1-z}.
\]
Thus $C^{-1}(z)$ is fixed by conjugation and lies in $F$.  Direct substitution shows that the formulas are inverse.  The final statement uses that $\No$ is real closed.
\end{proof}

\begin{proposition}[Cayley obstruction on mixed support]\label{prop:cayley-obstruction}
The Cayley transform acts on the full-support classical stratum but is not induced by a global split-projective linear automorphism of $\Pum^1(G(F[i]))$.
\end{proposition}

\begin{proof}
A matrix for the Cayley transform is
\[
\begin{pmatrix}
1&-i\\
1&i
\end{pmatrix}.
\]
Every entry is supported, so its Boolean support matrix is
\[
\begin{pmatrix}
1&1\\
1&1
\end{pmatrix},
\]
which is not invertible over $\B$.  Theorem \ref{thm:monomial-gl} therefore excludes invertibility over the split semiring.  Over the field $F[i]$, the determinant is nonzero and the classical action is valid on the full-support stratum.
\end{proof}

\subsection{Placing $\omega$ at a projective pole}

In the classical surcomplex projective line, the M\"obius transformation
\[
u=\frac{x}{\omega-x}
\]
is induced by
\[
[X:Y]\longmapsto[X:\omega Y-X].
\]
It sends $x=0$ to $u=0$ and $x=\omega$ to $u=\infty$.  This realizes $\omega$ as a projective pole rather than as a new arithmetic zero.

\begin{proposition}[The $\omega$-pole chart is not globally split-linear]\label{prop:omega-pole-obstruction}
The transformation $x\mapsto x/(\omega-x)$ is a classical projective automorphism of $\mathbf P^1(\No[i])$, but it does not extend as a global linear automorphism of $\Pum^1(G(\No[i]))$.
\end{proposition}

\begin{proof}
Its classical determinant is $\omega\neq0$, so it is a field-projective automorphism.  Any representing split matrix has a support pattern with at least three supported entries, hence its Boolean support matrix is not a permutation matrix.  Theorem \ref{thm:monomial-gl} excludes global split invertibility.
\end{proof}

\begin{remark}[Birthdays]
In Conway's construction, $0=\{\mid\}$ is born on day $0$, while $\omega^{-1}$ is born at a transfinite stage; see \cite{Conway,Simons}.  The external state $\tau$ is not an earlier surreal number.  One may assign it a separate support grade $\bot$ below all birthdays, but this is an enrichment of the birthday filtration, not a modification of Conway's recursive order.
\end{remark}

\section{Limits, poles, and removable indeterminacies}

The split formalism does not replace all limit theories by a single operation.  It separates a limit problem into an amplitude condition inside a monad and a support condition across the Boolean incidence skeleton.

\begin{definition}[Residue-monadic limit]\label{def:monadic-limit}
Let $q:A\twoheadrightarrow k$ be a split residue datum, let $a,L\in k$, choose a lift $\widetilde a\in A$, and let
\[
f:(\widetilde a+\mathfrak m)\setminus\{\widetilde a\}\longrightarrow A
\]
be a function.  We say that $f$ has \emph{residue-monadic limit} $L$ at $a$ if
\[
q(f(\widetilde a+h))=L
\]
for every nonzero $h\in\mathfrak m$ for which $f$ is defined.
\end{definition}

For finite hyperreals and an internal extension of a real function, this is the standard nonstandard criterion for a limit; see \cite{Robinson,Goldblatt}.  For the surcomplex valuation ring it states that all infinitesimal perturbations have the same complex standard part.

\begin{proposition}[Projective reciprocal extension]\label{prop:reciprocal-extension}
For every field $K$, the map
\[
\mathcal R:G(K)\longrightarrow\Pum^1(G(K)),
\qquad
x\longmapsto[1:x],
\]
extends ordinary reciprocal projectively.  It satisfies
\[
\mathcal R(h)=[h^{-1}:1]
\qquad(h\in K^\times),
\]
\[
\mathcal R(e)=\infty_e,
\qquad
\mathcal R(\tau)=\infty_\tau.
\]
\end{proposition}

\begin{proof}
If $h\in K^\times$, simultaneous multiplication by $h^{-1}$ gives
\[
[1:h]=[h^{-1}:1].
\]
The two zero cases are the definitions of $\infty_e$ and $\infty_\tau$.
\end{proof}

Thus a pole at supported zero and an absent denominator have distinct projective values.

\begin{theorem}[Resolution of a common-factor indeterminacy]\label{thm:common-factor-resolution}
Let $K$ be a field, let $a\in K$, and let $P,Q\in K[T]$.  Suppose
\[
P(T)=(T-a)^rP_0(T),
\qquad
Q(T)=(T-a)^rQ_0(T)
\]
for some $r\ge1$, with
\[
Q_0(a)\neq0.
\]
Then the rational projective map
\[
\Phi(x):=[P(x):Q(x)]
\]
agrees for $x\neq a$ with
\[
\widetilde\Phi(x):=[P_0(x):Q_0(x)],
\]
and $\widetilde\Phi$ extends to $x=a$ with value
\[
\widetilde\Phi(a)=[P_0(a):Q_0(a)].
\]
In $G(K)$, the unresolved pair at $x=a$ is
\[
[e:e],
\]
which has full support but is non-unimodular.  It is distinct from the absent pair
\[
[\tau:\tau].
\]
\end{theorem}

\begin{proof}
For $x\neq a$, the scalar $(x-a)^r$ is a nonzero unit of $K$, so
\[
[P(x):Q(x)]
=
[(x-a)^rP_0(x):(x-a)^rQ_0(x)]
=
[P_0(x):Q_0(x)].
\]
At $x=a$, the second coordinate $Q_0(a)$ is a unit, so $[P_0(a):Q_0(a)]$ is unimodular and defines a projective point.  Before cancellation, both $P(a)$ and $Q(a)$ are the supported ring zero $e$, giving the full-support but non-unimodular pair $(e,e)$.  The pair $(\tau,\tau)$ has empty support, so the two are distinct in $G(K)^2$.
\end{proof}

\begin{example}[The form $0/0$ at $x=1$]
For
\[
f(x)=\frac{x^2-1}{x-1},
\]
Theorem \ref{thm:common-factor-resolution} gives
\[
[x^2-1:x-1]=[x+1:1]
\]
away from $x=1$, and the resolved value is
\[
[2:1].
\]
If $h$ is a nonzero hyperreal or surcomplex infinitesimal, then
\[
f(1+h)=\frac{2h+h^2}{h}=2+h,
\]
so its standard part is $2$.  The Newtonian limit, the Leibniz infinitesimal computation, and the projective common-factor resolution therefore give the same amplitude result, while split support additionally distinguishes $(e,e)$ from $(\tau,\tau)$.
\end{example}

\begin{theorem}[Two-component split limit datum]\label{thm:two-component-limit}
For a map defined on a mixed-support punctured neighbourhood, a split limit consists of two independent conditions:
\begin{enumerate}[label=\textup{(\alph*)}]
\item constancy of the residue or standard part on each active infinitesimal fibre;
\item compatibility of the induced support map on the Boolean star of the base point.
\end{enumerate}
Neither condition implies the other.
\end{theorem}

\begin{proof}
The cubical decomposition of Theorem \ref{thm:cubical-monad} separates a neighbourhood into disjoint strata indexed by support activation masks.  Within each stratum, the residue map forgets infinitesimal amplitudes and tests condition (a).  The support character forgets all amplitudes and tests condition (b).  A function can preserve residues while changing a support mask, or preserve support while changing residues, so the conditions are independent.
\end{proof}

\section{Identity levels and univalence}

There are three mathematically distinct senses in which two infinitesimal objects may be called identical.

\begin{definition}[Three identity levels]\label{def:identity-levels}
Let $S=G(K)$.
\begin{enumerate}[label=\textup{(\roman*)}]
\item \emph{Arithmetic identity} is equality in $S$ or in a specified semiring quotient.
\item \emph{Support-modal identity} is equality after applying $\chi_K:S\to\B$.
\item \emph{Structural identity} is isomorphism or equivalence of pointed monads, formal neighbourhoods, or projective local geometries.
\end{enumerate}
\end{definition}

\begin{proposition}[Arithmetic and modal identities]\label{prop:arithmetic-modal}
For every $h\in K^\times$,
\[
e\neq h
\]
in $S$, but
\[
\chi_K(e)=\chi_K(h)=1.
\]
Freely adjoining the arithmetic identity $e=h$ yields $\B$, while adjoining it only on coordinates in $A$ yields
\[
\B^A\times S^{I\setminus A}.
\]
\end{proposition}

\begin{proof}
The inequality is set-theoretic because $e=0_K$ and $h\neq0$.  The support equality follows from the definition of $\chi_K$.  The quotient statements are Theorems \ref{thm:strict-identification-collapse} and \ref{thm:scoped-identity}.
\end{proof}

\begin{proposition}[Structural identity of zero and infinity monads]\label{prop:structural-identity}
Projective inversion gives an isomorphism of pointed split monads
\[
(\mu^\square(0_e),0_e)
\xrightarrow{\sim}
(\mu^\square(\infty_e),\infty_e)
\]
and likewise for the unsupported points $0_\tau$ and $\infty_\tau$.
\end{proposition}

\begin{proof}
This is Theorem \ref{thm:zero-infinity-monads}, with the distinguished centres included in the data.
\end{proof}

\begin{theorem}[Univalent interpretation]\label{thm:univalent-interpretation}
Assume a univalent universe containing the type of pointed split monads, with equivalences as its notion of structure-preserving equivalence.  Then the equivalence of Proposition \ref{prop:structural-identity} determines an identity path
\[
(\mu^\square(0_e),0_e)
=
(\mu^\square(\infty_e),\infty_e)
\]
in that universe.  This path does not induce an arithmetic equality
\[
e=\omega^{-1}
\]
in $G(\No[i])$.
\end{theorem}

\begin{proof}
The univalence axiom identifies the identity type between objects of a universe with the type of equivalences between them; see \cite{HoTT}.  Apply it to the projective inversion equivalence.  Arithmetic equality is a different identity type, namely equality of elements of the semiring.  By Theorem \ref{thm:strict-identification-collapse}, imposing $e=\omega^{-1}$ as an arithmetic relation produces the Boolean quotient and therefore cannot be obtained merely from the equivalence of pointed monads.
\end{proof}

\begin{remark}
Univalence therefore validates a genuine identity, but at the level at which the equivalence exists.  It does not erase invariants that belong to a stricter structure.  In particular,
\[
e^2=e,
\qquad
(\omega^{-1})^2=\omega^{-2}\neq\omega^{-1},
\]
so no multiplicative equivalence of pointed amplitude elements can identify them.
\end{remark}

\section{Final synthesis}

\begin{theorem}[Split-zero geometric synthesis]\label{thm:final-synthesis}
Let $R$ be a nonzero commutative ring, let $S=G(R)$, and let $e=0_R\neq\tau=0_S$.  The following statements hold.
\begin{enumerate}[label=\textup{(\arabic*)},leftmargin=3.2em]
\item The universal localization tower is
\[
G(R)
\longrightarrow
G(\Frac R)
\longrightarrow
\B
\longrightarrow
\zero
\]
when $R$ is a domain, where the successive operations invert classically nonzero amplitudes, the supported zero, and the actual semiring zero.

\item The absolute ring-null halo of $S^I$ is the Boolean cube $\B^I$.  Its elements are supported-zero masks; addition is union and multiplication is intersection.

\item Mixed localization at $x\in G(K)^I$ preserves full split arithmetic on nonzero-amplitude coordinates, Booleanizes supported-zero coordinates, and erases absent coordinates.

\item Adjoining an identity between a supported-zero mask and nonzero amplitudes Booleanizes exactly the selected support face and leaves the complementary coordinates unchanged.

\item The spectrum $\Spec G(R)$ is the one-point generic suspension of $\Spec R$.  The new point has Boolean stalk, while the old points have stalks $G(R_{\mathfrak p})$.

\item The exact $\mathbb F_1$ datum is the multiplicative monoid together with the preaddition inherited from $R$; the Boolean quotient is its support completion.

\item For a field $K$, the unimodular split projective space is the disjoint union of classical projective spaces over all nonempty Boolean faces:
\[
\Pum^n(G(K))
\cong
\bigsqcup_{\varnothing\neq A\subseteq[n]_0}
\mathbf P^{|A|-1}(K).
\]

\item Its finite-field point count and zeta function are given by Theorems \ref{thm:finite-field-point-count} and \ref{thm:support-zeta-product}.

\item Its global linear symmetry is monomial:
\[
\GL_{n+1}(G(K))
\cong
(K^\times)^{n+1}\rtimes S_{n+1}.
\]
Thus the globally valid projective symmetry is toric or $\mathbb F_1$-type.

\item For every residue map $A\twoheadrightarrow k$, the actual Leibniz monad is the supported fibre $a+\mathfrak m$.  Split zero augments it by a Boolean cube of support activations.  First-order neighbourhoods are Boolean cubes of affine spaces over $\mathfrak m/\mathfrak m^2$.

\item Formal completion commutes with split-zero globalization.

\item For finite hyperreals, the supported monads are the ordinary nonstandard monads.  For finite surcomplex numbers, they are the cosets $c+\mathfrak m_K$, with one complex graded line at each surreal valuation order.

\item The complete local surcomplex datum is
\[
(\text{support},\text{valuation order},\text{complex phase}).
\]
The element $e$ is the supported valuation-infinite centre, while $\omega^{-\alpha}$ is a nonzero invertible infinitesimal of finite order $\alpha$.

\item Projective inversion identifies the structured zero and infinity monads.  Under univalence this equivalence becomes an identity of structured objects, but strict arithmetic identification of $e$ with any nonzero amplitude collapses the corresponding coordinate to Boolean support.
\end{enumerate}
\end{theorem}

\begin{proof}
Items (1)--(14) are respectively Theorem \ref{thm:division-trichotomy}, Theorem \ref{thm:product-ring-reflection}, Theorem \ref{thm:mixed-localization}, Theorem \ref{thm:scoped-identity}, Theorem \ref{thm:generic-suspension} and Corollary \ref{cor:split-stalks}, Theorem \ref{thm:blueprint-presentation}, Theorem \ref{thm:projective-face-decomposition}, Theorems \ref{thm:finite-field-point-count} and \ref{thm:support-zeta-product}, Theorem \ref{thm:monomial-gl}, Theorems \ref{thm:monad-fibre}, \ref{thm:cubical-monad}, and \ref{thm:first-order-cubical}, Theorem \ref{thm:completion-commutes}, Corollaries \ref{cor:hyperreal-monads} and \ref{cor:surcomplex-monad}, Theorems \ref{thm:surcomplex-standard-part} and \ref{thm:complex-graded-layers} together with Definition \ref{def:support-order-phase}, and Theorems \ref{thm:zero-infinity-monads}, \ref{thm:strict-identification-collapse}, and \ref{thm:univalent-interpretation}.
\end{proof}

\begin{thebibliography}{99}

\bibitem{Globalization}
The Clankers,
\emph{A Note on the Globalization Semiring $G(\Z)$},
Zenodo, 2026.
DOI: \href{https://doi.org/10.5281/zenodo.18976982}{10.5281/zenodo.18976982}.

\bibitem{Secondary}
The Clankers,
\emph{A Secondary Note on Split-Zero Globalization: Semimodule Classification, Multiplicative Reconstruction, and All-Orders Zeta Deformation},
Zenodo, 2026.
DOI: \href{https://doi.org/10.5281/zenodo.19120905}{10.5281/zenodo.19120905}.

\bibitem{Homogenization}
The Clankers,
\emph{Support-Idempotent Homogenization for a Split-Zero Collatz F-Series Frame},
Zenodo, 2026.
DOI: \href{https://doi.org/10.5281/zenodo.19900461}{10.5281/zenodo.19900461}.

\bibitem{BerarducciMantova}
A.~Berarducci and V.~Mantova,
\emph{Surreal numbers, derivations and transseries},
J. Eur. Math. Soc. \textbf{20} (2018), no.~2, 339--390.
DOI: \href{https://doi.org/10.4171/JEMS/769}{10.4171/JEMS/769}.

\bibitem{Conway}
J.~H.~Conway,
\emph{On Numbers and Games},
2nd ed., A K Peters, Natick, MA, 2001.

\bibitem{Deitmar}
A.~Deitmar,
\emph{Schemes over $\mathbb F_1$},
in \emph{Number Fields and Function Fields---Two Parallel Worlds},
Progr. Math. \textbf{239}, Birkh\"auser, Boston, 2005, 87--100.

\bibitem{Golan}
J.~S.~Golan,
\emph{Semirings and Their Applications},
Kluwer Academic Publishers, Dordrecht, 1999.

\bibitem{Goldblatt}
R.~Goldblatt,
\emph{Lectures on the Hyperreals: An Introduction to Nonstandard Analysis},
Graduate Texts in Mathematics \textbf{188}, Springer, New York, 1998.

\bibitem{GrothendieckEGAII}
A.~Grothendieck,
\emph{\'El\'ements de g\'eom\'etrie alg\'ebrique II: \'Etude globale \'el\'ementaire de quelques classes de morphismes},
Publ. Math. IH\'ES \textbf{8} (1961), 5--222.

\bibitem{HoTT}
The Univalent Foundations Program,
\emph{Homotopy Type Theory: Univalent Foundations of Mathematics},
Institute for Advanced Study, 2013.
\url{https://homotopytypetheory.org/book/}.

\bibitem{Lorscheid}
O.~Lorscheid,
\emph{The geometry of blueprints. Part I: Algebraic background and scheme theory},
Adv. Math. \textbf{229} (2012), no.~3, 1804--1846.

\bibitem{Robinson}
A.~Robinson,
\emph{Non-standard Analysis},
North-Holland, Amsterdam, 1966.

\bibitem{Simons}
J.~Simons,
\emph{Meet the Surreal Numbers},
expository notes, 4 April 2017.

\end{thebibliography}

\end{document}
